\paragraph{像荷镜像能量与加权 Fekete 变分：Pick--Poisson 负谱体积势的上半平面实现}
本段把定义 \ref{def:xi-generalized-defect-entropy-logdet} 的广义缺陷熵
\(
S_{\mathrm{gen}}=-\log\det\mathsf P
\)
提升为上半平面 Dirichlet Green 能量的离散化，并由此给出平衡方程、稳定子群协变律与变分原理等一组可独立引用的静电学刻画。

\begin{proposition}[像荷镜像能量表示：\texorpdfstring{$-\log\det\mathsf P$}{-log det P} 等于上半平面 Dirichlet Green 能量的离散化]\label{prop:xi-pick-poisson-image-charge-green-energy}
\leanverified{paper\_xi\_pick\_poisson\_image\_charge\_green\_energy}
设有限点状缺陷点为 \(\{a_j\}_{j=1}^{\kappa}\subset\mathbb{D}\)（按重数展开），并令 \(\mathsf P\) 为其 Pick--Poisson Gram 矩阵 \eqref{eq:xi-pick-poisson-gram}。
定义“端点规一化”的上半平面坐标
\[
z_j:=\mathcal{C}^{-1}(a_j)\in\mathbb{H},
\qquad
t_j:=-\frac{1}{z_j}\in\mathbb{H},
\qquad
t_j=\gamma_j+\mathrm{i}y_j.
\]
则端点 Poisson 权重满足
\[
\boxed{\;p_j:=P_{a_j}(-1)=y_j=\Im(t_j).\;}
\]
并且伪双曲距离满足
\[
\boxed{\;
\rho(a_j,a_k)=\left|\frac{t_j-t_k}{t_j-\overline t_k}\right|\qquad(1\le j,k\le\kappa).\;}
\]
从而广义缺陷熵（负谱体积势）
\(
S_{\mathrm{gen}}:=-\log\det\mathsf P
\)
具有严格等价式
\[
\boxed{\;
S_{\mathrm{gen}}
=\sum_{j=1}^{\kappa}\bigl(-\log y_j\bigr)
\;+\;
2\sum_{1\le j<k\le\kappa}\Bigl(\log|t_j-\overline t_k|-\log|t_j-t_k|\Bigr).\;}
\]
令上半平面 Dirichlet Green 函数为
\[
G_{\mathbb{H}}(t,s):=\log\left|\frac{t-\overline s}{t-s}\right|,
\]
则上式可写为
\[
S_{\mathrm{gen}}
=\sum_{j=1}^{\kappa}\bigl(-\log\Im(t_j)\bigr)
\;+\;
2\sum_{1\le j<k\le\kappa}G_{\mathbb{H}}(t_j,t_k).
\]
因此点状缺陷在端点视界下的“压缩能量”具有严格静电学实现：外部能量完全由真实点与其边界镜像的 Green 相互作用封闭，而自能项 \(-\sum_j\log\Im(t_j)\) 给出端点吸收的不可消去代价。
\end{proposition}
\begin{proof}[证明要点]
由命题 \ref{prop:xi-pick-poisson-det-scaling-law} 的计算，
\(
P_w(-1)=\Im z/|z|^2
\)（其中 \(z=\mathcal{C}^{-1}(w)\)）。
取 \(w=a_j\) 并用 \(\Im(-1/z)=\Im z/|z|^2\) 得 \(p_j=\Im(t_j)=y_j\)。
\par
伪双曲距离在 Cayley 变换下满足 \(\rho(a_j,a_k)=\left|\frac{z_j-z_k}{z_j-\overline z_k}\right|\)，而映射 \(z\mapsto -1/z\) 属于 \(\mathrm{Aut}(\mathbb{H})\) 并保持该不变量，故亦有 \(\rho(a_j,a_k)=\left|\frac{t_j-t_k}{t_j-\overline t_k}\right|\)。
将 \(p_j\) 与 \(\rho_{jk}\) 代入命题 \ref{prop:xi-pick-poisson-det-factorization} 的分解并取负对数，即得能量表示。证毕。
\end{proof}

\begin{theorem}[端点锚定伪双曲几何的横向分离反演公式]\label{thm:xi-pick-poisson-anchored-transverse-separation-inversion}
\leanverified{paper\_xi\_pick\_poisson\_anchored\_transverse\_separation\_inversion}
沿用命题 \ref{prop:xi-pick-poisson-image-charge-green-energy} 的设定。
对缺陷点 \(a_j,a_k\in\mathbb D\) 令端点 Poisson 权重
\(
p_j:=P_{a_j}(-1),\ p_k:=P_{a_k}(-1)
\)，
并令端点规一化上半平面坐标
\(
t_j=\gamma_j+\mathrm{i}p_j,\ t_k=\gamma_k+\mathrm{i}p_k\in\mathbb H
\)。
记 \(\rho_{jk}:=\rho(a_j,a_k)\)。
则对任意 \(j\neq k\)，横向分离
\[
\Delta_{jk}:=|\gamma_j-\gamma_k|
\]
可由 \((p_j,p_k,\rho_{jk})\) 显式反演为
\[
\boxed{\;
\Delta_{jk}^2
=\frac{\rho_{jk}^2(p_j+p_k)^2-(p_j-p_k)^2}{1-\rho_{jk}^2}\; }.
\]
并且可行性约束（右端非负）等价于
\[
\rho_{jk}\ \ge\ \frac{|p_j-p_k|}{p_j+p_k},
\]
其中等号当且仅当 \(\Delta_{jk}=0\)。
\end{theorem}
\begin{proof}
由命题 \ref{prop:xi-pick-poisson-image-charge-green-energy} 的半平面表示，
\[
\rho_{jk}
=\left|\frac{t_j-t_k}{t_j-\overline t_k}\right|.
\]
写
\(
t_j-t_k=(\gamma_j-\gamma_k)+\mathrm{i}(p_j-p_k)
\)
与
\(
t_j-\overline t_k=(\gamma_j-\gamma_k)+\mathrm{i}(p_j+p_k)
\)，
即得
\[
\rho_{jk}^2
=\frac{(\gamma_j-\gamma_k)^2+(p_j-p_k)^2}{(\gamma_j-\gamma_k)^2+(p_j+p_k)^2}
=\frac{\Delta_{jk}^2+(p_j-p_k)^2}{\Delta_{jk}^2+(p_j+p_k)^2}.
\]
移项解 \(\Delta_{jk}^2\) 得到盒中公式。
右端非负等价于 \(\rho_{jk}^2(p_j+p_k)^2\ge (p_j-p_k)^2\)，亦即所示不等式；
当且仅当取等号时 \(\Delta_{jk}=0\)。证毕。
\end{proof}

\begin{theorem}[Pick--Poisson Gram 数据的几何刚性：由 \texorpdfstring{$(p_j,\rho_{jk})$}{(p_j,rho_jk)} 复原锚定缺陷配置]\label{thm:xi-pick-poisson-anchored-gram-rigidity-reconstruction}
取 \(\kappa\ge 2\)。
设给定端点 Poisson 权重 \(p_j>0\)（\(1\le j\le\kappa\)）与两两伪双曲距离 \(\rho_{jk}\in(0,1)\)（\(1\le j<k\le\kappa\)），且它们可由某个缺陷多重集 \(A=\{a_j\}_{j=1}^{\kappa}\subset\mathbb D\) 实现。
沿用命题 \ref{prop:xi-pick-poisson-image-charge-green-energy} 的端点规一化上半平面坐标 \(t_j=\gamma_j+\mathrm{i}p_j\)。
则：
\begin{enumerate}
  \item 由定理 \ref{thm:xi-pick-poisson-anchored-transverse-separation-inversion} 可逐对计算横向距离
  \[
  d_{jk}:=\Delta_{jk}=|\gamma_j-\gamma_k|\qquad(1\le j<k\le\kappa).
  \]
  \item 距离数据 \(\{d_{jk}\}_{j<k}\) 唯一决定点列 \(\{\gamma_j\}\subset\RR\) 在整体刚体对称（平移与反射）意义下的同构类。
  \item 从而 \(\{t_j\}\subset\mathbb H\) 在同一平移与反射对称下确定；进而通过
  \[
  z_j:=-\frac{1}{t_j},\qquad a_j:=\mathcal C(z_j)
  \]
  （命题 \ref{prop:xi-pick-poisson-image-charge-green-energy} 的记号）复原缺陷点多重集 \(A\) 在端点稳定子群所诱导的等距对称下的同构类。
\end{enumerate}
并且存在显式复原流程：任选规范化 \(\gamma_1=0\)。
若存在某个 \(r\in\{2,\dots,\kappa\}\) 使 \(d_{1r}>0\)，则固定取向为 \(\gamma_r=d_{1r}\ge 0\)；
对每个 \(k\notin\{1,r\}\)，由二方程组
\[
|\gamma_k-\gamma_1|=d_{1k},\qquad |\gamma_k-\gamma_r|=d_{rk}
\]
在既定取向下唯一确定 \(\gamma_k\)；最终取 \(t_j=\gamma_j+\mathrm{i}p_j\) 并按上式得到 \(a_j\)。
若对所有 \(r\ge 2\) 都有 \(d_{1r}=0\)，则 \(\gamma_1=\cdots=\gamma_\kappa=0\)。
\end{theorem}
\begin{proof}
第 (1) 点由定理 \ref{thm:xi-pick-poisson-anchored-transverse-separation-inversion} 逐对计算 \(\Delta_{jk}\) 立得。
\par
对第 (2) 点，这是 \(\RR\) 上的距离几何刚性：已知全部 \(d_{jk}=|\gamma_j-\gamma_k|\) 后，先以 \(\gamma_1=0\) 消去平移自由度。
此时每个 \(\gamma_k\) 只能取 \(\pm d_{1k}\)。
若进一步规定 \(\gamma_2=d_{12}\ge 0\) 固定取向，则对任意 \(k\ge 3\)，约束 \(|\gamma_k-\gamma_2|=d_{2k}\) 排除其中一个符号分支，故 \(\gamma_k\) 在该规范化下唯一确定。
由此得到的 \(\{\gamma_j\}\) 与任何其他实现仅相差整体平移或反射。
\par
第 (3) 点由 \(t_j\mapsto z_j:=-1/t_j\mapsto a_j=\mathcal C(z_j)\) 的双射性推出，并且反射 \(\gamma\mapsto-\gamma\) 与平移 \(\gamma\mapsto\gamma+c\) 在 \(\mathbb H\) 上诱导的等距对称经 Cayley 变换传递为 \(\mathbb D\) 的端点稳定子群对称。
显式构造流程即为上述规范化下的逐点解算。证毕。
\end{proof}

\begin{remark}[高度与横向分离的完备性]\label{rem:xi-pick-poisson-anchored-gram-rigidity-sufficient-statistic}
定理 \ref{thm:xi-pick-poisson-anchored-gram-rigidity-reconstruction} 表明：
端点 Poisson 权重 \(\{p_j\}\) 与两两伪双曲距离 \(\{\rho_{jk}\}\) 共同构成端点锚定模型中缺陷多重集的几何完备统计量；
其中 \(\{p_j\}\) 在上半平面实现中给出高度坐标，而 \(\{\rho_{jk}\}\) 通过定理 \ref{thm:xi-pick-poisson-anchored-transverse-separation-inversion} 决定横向分离矩阵，从而闭合地恢复全体相对构型信息。
\end{remark}

\endinput

% Pick--Poisson determinant: Vandermonde / resultant / discriminant identities.

\begin{theorem}[Pick--Poisson 行列式的 Vandermonde--结果式恒等式]\label{thm:xi-pick-poisson-vandermonde-resultant-identity}
取互异点集 \(A=\{a_1,\dots,a_\kappa\}\subset\mathbb D\)，并令 Pick--Poisson Gram 矩阵 \(\mathsf P\) 由 \eqref{eq:xi-pick-poisson-gram} 给出。
则 \(\det\mathsf P>0\)，并满足乘积分解（命题 \ref{prop:xi-pick-poisson-det-factorization} 的记号）
\[
\det\mathsf P
=
\Bigl(\prod_{j=1}^{\kappa} \frac{1-|a_j|^2}{|1+a_j|^2}\Bigr)\cdot
\Bigl(\prod_{1\le j<k\le \kappa}\rho(a_j,a_k)^2\Bigr),
\qquad
\rho(a,b):=\left|\frac{a-b}{1-\overline b a}\right|.
\]
进一步，令
\[
F(z):=\prod_{j=1}^\kappa (z-a_j)\quad(\text{首一多项式}),\qquad
F^\ast(z):=z^\kappa\,\overline{F(1/\overline z)}=\prod_{j=1}^\kappa (1-\overline a_j z).
\]
则成立纯代数不变量恒等式
\[
\det\mathsf P
=
\frac{\displaystyle \prod_{1\le j<k\le\kappa}|a_j-a_k|^2\ \cdot\ \prod_{j=1}^\kappa (1-|a_j|^2)^2}
{\displaystyle |\Res(F,F^\ast)|\ \cdot\ \prod_{j=1}^\kappa |1+a_j|^2},
\]
以及平方化形式
\[
(\det\mathsf P)^2
=
\frac{\displaystyle |\Disc(F)|\ \cdot\ \prod_{j=1}^\kappa (1-|a_j|^2)^4}
{\displaystyle |\Res(F,F^\ast)|^2\ \cdot\ \prod_{j=1}^\kappa |1+a_j|^4}.
\]
\leanverified{paper\_zeta\_vandermonde\_resultant\_seeds}
\leanverified{paper\_zeta\_vandermonde\_resultant\_package}
\leanverified{paper\_xi\_pick\_poisson\_vandermonde\_resultant\_identity}
\end{theorem}
\begin{proof}
由定义
\[
\prod_{1\le j<k\le\kappa}\rho(a_j,a_k)^2
=
\prod_{1\le j<k\le\kappa}\frac{|a_j-a_k|^2}{|1-\overline a_k a_j|^2}.
\]
另一方面，结果式的根值表述给出
\[
\Res(F,F^\ast)=\prod_{j=1}^\kappa F^\ast(a_j)
=\prod_{j=1}^\kappa\ \prod_{k=1}^\kappa (1-\overline a_k a_j).
\]
取模得
\[
|\Res(F,F^\ast)|
=
\Bigl(\prod_{j=1}^\kappa (1-|a_j|^2)\Bigr)\cdot
\Bigl(\prod_{1\le j<k\le\kappa}|1-\overline a_k a_j|^2\Bigr).
\]
将该式与 \(\rho\) 的展开代入 \(\det\mathsf P\) 的乘积分解并整理，即得第一恒等式。
最后，由首一多项式判别式满足
\(
|\Disc(F)|=\prod_{1\le j<k\le\kappa}|a_j-a_k|^2\cdot \prod_{1\le j<k\le\kappa}|a_j-a_k|^2
=\prod_{j<k}|a_j-a_k|^4
\)，
从而平方化得到第二恒等式。证毕。
\end{proof}

\begin{corollary}[行列式下界推出伪双曲分离的充分条件]\label{cor:xi-pick-poisson-determinant-separation-sufficient}
\leanverified{paper\_xi\_pick\_poisson\_determinant\_separation\_sufficient}
在定理 \ref{thm:xi-pick-poisson-vandermonde-resultant-identity} 的设定下，若进一步满足 \(\max_j|a_j|\le r<1\)，则
\[
\min_{1\le j<k\le\kappa}\rho(a_j,a_k)
\ \ge\
(\det\mathsf P)^{1/2}\,(1-r)^{\kappa}.
\]
\end{corollary}
\begin{proof}
由 \(|a_j|\le r\) 得 \(|1+a_j|\ge 1-|a_j|\ge 1-r\)，故
\(
\frac{1-|a_j|^2}{|1+a_j|^2}\le (1-r)^{-2}
\)。
代入 \(\det\mathsf P=\bigl(\prod_j \frac{1-|a_j|^2}{|1+a_j|^2}\bigr)\bigl(\prod_{j<k}\rho_{jk}^2\bigr)\) 得
\[
\prod_{1\le j<k\le\kappa}\rho(a_j,a_k)^2
\ \ge\
\det\mathsf P\cdot (1-r)^{2\kappa}.
\]
由于每个 \(\rho(a_j,a_k)^2\in(0,1]\)，故其乘积不超过任何单个因子，特别地
\[
\prod_{j<k}\rho(a_j,a_k)^2
\ \le\
\min_{j<k}\rho(a_j,a_k)^2.
\]
合并并取平方根即得结论。证毕。
\end{proof}

\endinput


\begin{proposition}[平衡方程：\texorpdfstring{$S_{\mathrm{gen}}$}{S_gen} 的临界点满足镜像库仑力严格抵消]\label{prop:xi-pick-poisson-energy-balance-equations}
\leanverified{paper\_xi\_pick\_poisson\_energy\_balance\_equations}
在命题 \ref{prop:xi-pick-poisson-image-charge-green-energy} 的上半平面参数 \((\gamma_j,y_j)\) 下，将 \(S_{\mathrm{gen}}\) 视为变量函数。
则其欧氏梯度具有显式“镜像力差”形式：对任意 \(j\),
\[
\boxed{\;
\frac{\partial S_{\mathrm{gen}}}{\partial \gamma_j}
=
2\sum_{k\neq j}\left(
\frac{\gamma_j-\gamma_k}{(\gamma_j-\gamma_k)^2+(y_j+y_k)^2}
-
\frac{\gamma_j-\gamma_k}{(\gamma_j-\gamma_k)^2+(y_j-y_k)^2}
\right)\;}
\]
并且
\[
\boxed{\;
\frac{\partial S_{\mathrm{gen}}}{\partial y_j}
=
-\frac{1}{y_j}
+
2\sum_{k\neq j}\left(
\frac{y_j+y_k}{(\gamma_j-\gamma_k)^2+(y_j+y_k)^2}
-
\frac{y_j-y_k}{(\gamma_j-\gamma_k)^2+(y_j-y_k)^2}
\right).\;}
\]
因此任何内部极值配置必须满足上述 \(\kappa\) 组方程同时为零；其含义为：每个点在边界导体（镜像）作用下的切向与法向库仑力均达到严格平衡，而自能项 \(-\log y_j\) 仅在法向方程中提供不可消去的端点吸引。
\end{proposition}
\begin{proof}[证明要点]
由命题 \ref{prop:xi-pick-poisson-image-charge-green-energy} 的能量展开，
\[
S_{\mathrm{gen}}
=-\sum_j\log y_j
\;+\;\sum_{j<k}\log\!\bigl((\gamma_j-\gamma_k)^2+(y_j+y_k)^2\bigr)
\;-\;\sum_{j<k}\log\!\bigl((\gamma_j-\gamma_k)^2+(y_j-y_k)^2\bigr).
\]
对 \(\gamma_j\)、\(y_j\) 分别逐项求偏导并整理对称项即可得到所示表达式。证毕。
\end{proof}

\begin{corollary}[视界稳定子群协变律：平移不变、缩放具 \texorpdfstring{$\kappa\log\lambda$}{k log lambda} 余因子]\label{cor:xi-pick-poisson-energy-stabilizer-covariance}
\leanverified{paper\_xi\_pick\_poisson\_energy\_stabilizer\_covariance}
令上半平面点集作水平平移 \(t_j\mapsto t_j+x_0\)（\(x_0\in\RR\)）与正缩放 \(t_j\mapsto \lambda t_j\)（\(\lambda>0\)）。
则
\[
\boxed{\;
S_{\mathrm{gen}}(t_1+x_0,\dots,t_\kappa+x_0)=S_{\mathrm{gen}}(t_1,\dots,t_\kappa),\;}
\]
并且
\[
\boxed{\;
S_{\mathrm{gen}}(\lambda t_1,\dots,\lambda t_\kappa)=S_{\mathrm{gen}}(t_1,\dots,t_\kappa)-\kappa\log\lambda.\;}
\]
换言之，点间互能项（Green 差）在视界稳定子群下严格不变，而自能项通过 \(-\sum_j\log\Im(t_j)\) 产生唯一可计算的同调余因子 \(-\kappa\log\lambda\)。
\end{corollary}
\begin{proof}
平移保持各差分 \(t_j-t_k\)、\(t_j-\overline t_k\) 且保持 \(\Im(t_j)\)，故能量不变。
缩放下 \(\Im(\lambda t_j)=\lambda\Im(t_j)\)，而 \(\log|\lambda(t_j-\overline t_k)|-\log|\lambda(t_j-t_k)|\) 的 \(\log\lambda\) 项相消，故仅自能项贡献 \(-\kappa\log\lambda\)。证毕。
\end{proof}

% (User input window)

\begin{corollary}[缺陷度作为对数尺度的熵斜率：\texorpdfstring{$\frac{\dd}{\dd(\log m)}S_{\mathrm{gen}}(m)=\kappa$}{d/dlogm S_gen(m)=kappa}]\label{cor:xi-pick-poisson-entropy-derivative-visible-dimension}
\leanpartial{paper\_xi\_pick\_poisson\_entropy\_derivative\_visible\_dimension}{the requested Lean signature is false without a positivity hypothesis on \texttt{detP 1}; taking \texttt{detP = 0} and \texttt{Sgen = 0} satisfies the assumptions but not the claimed affine law}
沿用推论 \ref{cor:xi-pick-poisson-energy-stabilizer-covariance} 的记号，并固定一组上半平面点 \(t_1,\dots,t_\kappa\in\mathbb{H}\)。
对任意尺度参数 \(m>0\) 定义缩放后的点集 \(t^{(m)}_j:=t_j/m\)，并令 \(\mathsf P(m):=\mathsf P(t^{(m)}_1,\dots,t^{(m)}_\kappa)\) 为相应 Pick--Poisson Gram 矩阵（\eqref{eq:xi-pick-poisson-gram}）。
定义广义缺陷熵
\[
S_{\mathrm{gen}}(m):=-\log\det\mathsf P(m).
\]
则成立严格协变律
\[
\boxed{\;
\det\mathsf P(m)=m^{-\kappa}\det\mathsf P(1),\qquad
S_{\mathrm{gen}}(m)=S_{\mathrm{gen}}(1)+\kappa\log m\;}
\]
从而
\[
\boxed{\;\frac{\dd}{\dd(\log m)}S_{\mathrm{gen}}(m)=\kappa.\;}
\]
因此在 Pick--Poisson 口径下，缺陷度 \(\kappa\) 可被刻画为“边界体积熵”在对数尺度方向上的严格斜率：它不依赖几何细节，仅由外部化维数（缺陷度）决定。
\end{corollary}
\begin{proof}
由推论 \ref{cor:xi-pick-poisson-energy-stabilizer-covariance} 的缩放律，取 \(\lambda=1/m\) 得
\[
S_{\mathrm{gen}}(t_1/m,\dots,t_\kappa/m)=S_{\mathrm{gen}}(t_1,\dots,t_\kappa)+\kappa\log m.
\]
将左端记为 \(S_{\mathrm{gen}}(m)\) 即得所示线性律。
对两边关于 \(\log m\) 求导即得导数恒等式。证毕。
\end{proof}


\begin{proposition}[加权 Fekete 变分原理：最大化 \texorpdfstring{$\det\mathsf P$}{det P} 等价于最小化镜像能量]\label{prop:xi-pick-poisson-weighted-fekete-variational}
固定 \(\kappa\) 并取约束集合 \(\mathcal{K}\subset\mathbb{H}^{\kappa}\)。
将 \(S_{\mathrm{gen}}=-\log\det\mathsf P\) 理解为扩展实值函数：若 \(\det\mathsf P=0\) 则取 \(S_{\mathrm{gen}}=+\infty\)。
则极值问题
\[
\max_{(t_1,\dots,t_\kappa)\in\mathcal{K}}\ \det\mathsf P(t_1,\dots,t_\kappa)
\qquad\Longleftrightarrow\qquad
\min_{(t_1,\dots,t_\kappa)\in\mathcal{K}}\ S_{\mathrm{gen}}(t_1,\dots,t_\kappa)
\]
严格等价。
当 \(\mathcal{K}\) 紧致且包含至少一个使 \(\det\mathsf P>0\) 的配置时，上述最大值与最小值均存在且有限。
若进一步施加几何分离下界 \(\rho_{\min}\ge\rho_0>0\) 与高度下界 \(y_{\min}>0\)，则极值序列预紧，且任何内部极值点必满足命题 \ref{prop:xi-pick-poisson-energy-balance-equations} 的平衡方程。
\end{proposition}
\begin{proof}[证明要点]
等价性来自对数的单调性：
\(
\arg\max\det\mathsf P=\arg\min(-\log\det\mathsf P)
\)。
若 \(\mathcal{K}\) 紧致，则 \(\det\mathsf P\) 为连续函数，故最大值存在。
另一方面，若 \(\mathcal{K}\) 中存在一点使 \(\det\mathsf P>0\)，则 \(\inf_{\mathcal{K}}S_{\mathrm{gen}}<\infty\)；
并且 \(S_{\mathrm{gen}}\) 作为 \(-\log\circ\det\) 的扩展实值函数在 \(\mathcal{K}\) 上下半连续，故最小值存在且有限。
\par
分离/高度下界排除了贴边与碰撞退化（从而排除能量发散并保证 \(\det\mathsf P\) 在该域内一致远离零），给出预紧性；内部极值点满足梯度为零即得平衡方程。证毕。
\end{proof}

\begin{proposition}[行列式链式分解：条件通量的逐步记账律]\label{prop:xi-pick-poisson-schur-flux-ledger}
\leanverified{paper\_xi\_pick\_poisson\_schur\_flux\_ledger}
对任意固定点序 \(1,\dots,\kappa\)，记 \(\mathsf P^{(m)}\) 为前 \(m\) 个点诱导的主子矩阵。
定义 Schur 补（条件通量）
\[
\pi_m
:=
\mathsf P_{mm}-\mathsf P_{m,1:m-1}\,(\mathsf P^{(m-1)})^{-1}\,\mathsf P_{1:m-1,m},
\qquad
m=1,\dots,\kappa,
\]
其中约定 \(\mathsf P^{(0)}\) 为空。
则有严格乘积律
\[
\boxed{\;
\det\mathsf P=\prod_{m=1}^{\kappa}\pi_m,
\qquad
S_{\mathrm{gen}}=-\sum_{m=1}^{\kappa}\log\pi_m.\;}
\]
并且每一 \(\pi_m\) 满足
\[
\boxed{\;
0<\pi_m\le \mathsf P_{mm}=p_m.\;}
\]
因此 \(S_{\mathrm{gen}}\) 具有严格的“逐步外部化记账”结构：每次加入一个缺陷，新增能量恰等于其在既有缺陷张成空间上的条件通量损失。
\end{proposition}
\begin{proof}[证明要点]
\(\det\mathsf P=\prod_m\pi_m\) 为正定矩阵的 Schur 补链式分解（等价 Cholesky 分解）标准结论。
上界 \(\pi_m\le \mathsf P_{mm}\) 可由 Gram--Schmidt 解释：\(\mathsf P\) 为一族向量的 Gram 矩阵时，\(\pi_m\) 等于第 \(m\) 个向量在此前张成子空间的正交补上的平方范数，故不超过其全范数平方 \(\mathsf P_{mm}\)。证毕。
\end{proof}

\begin{corollary}[Schur 补通量的闭式：条件体积增量的显式乘积公式]\label{cor:xi-pick-poisson-schur-flux-closed-form}
沿用命题 \ref{prop:xi-pick-poisson-schur-flux-ledger} 的设定，并固定点序 \((a_1,\dots,a_\kappa)\)。
记 \(\mathsf P^{(m)}\) 为前 \(m\) 个点诱导的主子矩阵，并记 \(\pi_m\) 为相应 Schur 补通量（命题 \ref{prop:xi-pick-poisson-schur-flux-ledger}）；
则对任意 \(m\ge 2\) 有
\[
\boxed{\;
\pi_m=\frac{\det\mathsf P^{(m)}}{\det\mathsf P^{(m-1)}}
=p_m\prod_{j<m}\rho(a_j,a_m)^2\;}
\qquad\Bigl(p_m:=P_{a_m}(-1)\Bigr).
\]
在命题 \ref{prop:xi-pick-poisson-image-charge-green-energy} 的上半平面坐标 \(t_j=\gamma_j+\mathrm{i}p_j\) 下，上式等价为
\[
\pi_m
=p_m\prod_{j<m}\frac{|t_m-t_j|^2}{|t_m-\overline t_j|^2}
=p_m\prod_{j<m}
\frac{(\gamma_m-\gamma_j)^2+(p_m-p_j)^2}{(\gamma_m-\gamma_j)^2+(p_m+p_j)^2}.
\]
从而得到逐步记账恒等式
\[
-\log\pi_m
=-\log p_m
+2\sum_{j<m}\bigl(-\log\rho(a_j,a_m)\bigr).
\]
\end{corollary}
\begin{proof}
由命题 \ref{prop:xi-pick-poisson-schur-flux-ledger} 的 Schur 链式分解，\(\pi_m=\det\mathsf P^{(m)}/\det\mathsf P^{(m-1)}\)。
另一方面，对任意有限点集的 Pick--Poisson Gram 矩阵均有行列式分解（命题 \ref{prop:xi-pick-poisson-det-factorization}；亦见定理 \ref{thm:xi-pick-poisson-vandermonde-resultant-identity}）：
\[
\det\mathsf P^{(m)}
=\Bigl(\prod_{r\le m}p_r\Bigr)\Bigl(\prod_{1\le j<k\le m}\rho(a_j,a_k)^2\Bigr).
\]
将该式对 \(m\) 与 \(m-1\) 取比值即得
\(
\pi_m=p_m\prod_{j<m}\rho(a_j,a_m)^2
\)。
上半平面表达式由 \(\rho(a_j,a_m)=\left|\frac{t_m-t_j}{t_m-\overline t_j}\right|\)（命题 \ref{prop:xi-pick-poisson-image-charge-green-energy}）直接代入。
对数恒等式对乘积式取负对数即可。证毕。
\end{proof}

\endinput


\begin{corollary}[负谱体积势的近碰撞证书：行列式强制伪双曲收缩]\label{cor:xi-pick-poisson-near-collision-certificate}
\leanverified{paper\_xi\_pick\_poisson\_near\_collision\_certificate}
设 \(\kappa\ge 2\)，并令 \(\mathsf P\) 为缺陷点 \(\{a_j\}_{j=1}^{\kappa}\subset\mathbb{D}\) 的 Pick--Poisson Gram 矩阵。
记端点 Poisson 权重 \(p_j:=P_{a_j}(-1)\) 与伪双曲距离 \(\rho_{jk}:=\rho(a_j,a_k)\)。
定义交叉互能
\[
E_{\mathrm{int}}
:=\bigl(-\log\det\mathsf P\bigr)\ -\ \sum_{j=1}^{\kappa}\bigl(-\log p_j\bigr)
=2\sum_{1\le j<k\le\kappa}\bigl(-\log\rho_{jk}\bigr)\ \ge\ 0.
\]
则存在某对 \(j<k\) 使
\[
\rho_{jk}\ \le\
\left(\frac{\det\mathsf P}{\prod_{j=1}^{\kappa}p_j}\right)^{\!\frac1{\kappa(\kappa-1)}}
\ =\
\exp\!\left(-\frac{E_{\mathrm{int}}}{\kappa(\kappa-1)}\right).
\]
\end{corollary}
\begin{proof}
由命题 \ref{prop:xi-pick-poisson-det-factorization}，
\(
\det\mathsf P=\bigl(\prod_j p_j\bigr)\bigl(\prod_{j<k}\rho_{jk}^2\bigr)
\)。
记 \(\rho_{\min}:=\min_{j<k}\rho_{jk}\)，则 \(\prod_{j<k}\rho_{jk}^2\ge \rho_{\min}^{\kappa(\kappa-1)}\)，从而
\[
\rho_{\min}\le \left(\frac{\det\mathsf P}{\prod_j p_j}\right)^{1/(\kappa(\kappa-1))}.
\]
取到极小值的一对 \((j,k)\) 即得第一式。
第二式由 \(E_{\mathrm{int}}=-\log\!\bigl(\det\mathsf P/\prod_j p_j\bigr)\) 立得。证毕。
\end{proof}

\begin{corollary}[谱隙迫使几何分离：\texorpdfstring{$\rho_{\min}$}{rho\_min} 的下界与条件数二次下界]\label{cor:xi-pick-poisson-min-separation-lowerbound}
\leanverified{paper\_xi\_pick\_poisson\_min\_separation\_lowerbound}
沿用上述记号，并记
\[
p_\Sigma:=\sum_{j=1}^{\kappa}p_j=\mathrm{tr}(\mathsf P),\qquad
p_{\max}:=\max_{1\le j\le \kappa}p_j,\qquad
\rho_{\min}:=\min_{1\le j<k\le \kappa}\rho_{jk}.
\]
则成立严格下界
\[
\boxed{\;
\rho_{\min}^{2}\ \ge\ \frac{\det\mathsf P}{p_{\max}\,p_\Sigma^{\kappa-1}}\; }.
\]
并且 Pick--Poisson Gram 的条件数满足普适二次下界
\[
\boxed{\;
\kappa(\mathsf P):=\frac{\lambda_{\max}(\mathsf P)}{\lambda_{\min}(\mathsf P)}\ \ge\ \rho_{\min}^{-2}\; } .
\]
\end{corollary}
\begin{proof}
取到 \(\rho_{\min}\) 的指标对 \(j^\star<k^\star\)，并记 \(I=\{j^\star,k^\star\}\)。
由主子矩阵特征值交错性，
\(
\lambda_{\min}(\mathsf P)\le \lambda_{\min}(\mathsf P_I)
\)。
对 \(2\times 2\) 正定矩阵有
\(
\lambda_{\min}(\mathsf P_I)\le \det(\mathsf P_I)/\lambda_{\max}(\mathsf P_I)
\le \det(\mathsf P_I)/p_{j^\star}
\)；
并由 \(\det(\mathsf P_I)=p_{j^\star}p_{k^\star}\rho_{\min}^{2}\)（命题 \ref{prop:xi-pick-poisson-det-factorization} 在 \(\kappa=2\) 的特例）得
\[
\lambda_{\min}(\mathsf P)\ \le\ p_{k^\star}\rho_{\min}^{2}\ \le\ p_{\max}\rho_{\min}^{2}.
\]
另一方面，对任意正定矩阵有
\(
\lambda_{\min}(\mathsf P)\ge \det(\mathsf P)/(\mathrm{tr}\,\mathsf P)^{\kappa-1}=\det(\mathsf P)/p_\Sigma^{\kappa-1}
\)。
合并即得 \(\rho_{\min}^{2}\ge \det\mathsf P/(p_{\max}p_\Sigma^{\kappa-1})\)。
最后，由 \(\lambda_{\max}(\mathsf P)\ge p_{\max}\) 与上式 \(\lambda_{\min}(\mathsf P)\le p_{\max}\rho_{\min}^{2}\) 得
\(
\kappa(\mathsf P)=\lambda_{\max}/\lambda_{\min}\ge \rho_{\min}^{-2}
\)。
证毕。
\end{proof}

\begin{corollary}[Schur 补的单步坍缩下界：互能必须在某一步骤中集中显现]\label{cor:xi-pick-poisson-schur-one-step-collapse}
沿用命题 \ref{prop:xi-pick-poisson-schur-flux-ledger} 的记号。
令 \((\pi_m)_{m=1}^{\kappa}\) 为 Schur 补通量，使得 \(\det\mathsf P=\prod_{m=1}^{\kappa}\pi_m\) 且 \(0<\pi_m\le p_m\)。
则存在某个 \(m_\star\in\{1,\dots,\kappa\}\) 使
\[
\pi_{m_\star}\ \le\ p_{m_\star}\cdot
\left(\frac{\det\mathsf P}{\prod_{j=1}^{\kappa}p_j}\right)^{\!\frac1\kappa}
\ =\
p_{m_\star}\cdot\left(\prod_{1\le j<k\le\kappa}\rho_{jk}^2\right)^{\!\frac1\kappa}.
\]
等价地，
\[
-\log\pi_{m_\star}\ \ge\ -\log p_{m_\star}\ +\ \frac1\kappa\,E_{\mathrm{int}}.
\]
\leanverified{paper\_xi\_pick\_poisson\_schur\_one\_step\_collapse}
\end{corollary}
\begin{proof}
由 \(\det\mathsf P=\prod_m \pi_m\) 与命题 \ref{prop:xi-pick-poisson-det-factorization} 得
\[
\prod_{m=1}^{\kappa}\frac{p_m}{\pi_m}
=\frac{\prod_m p_m}{\det\mathsf P}
=\left(\prod_{j<k}\rho_{jk}^2\right)^{-1}.
\]
令 \(\eta_m:=p_m/\pi_m\ge 1\)。由 \(\max_m\eta_m\ge (\prod_m\eta_m)^{1/\kappa}\) 知存在 \(m_\star\) 使
\(
\eta_{m_\star}\ge \left(\prod_{j<k}\rho_{jk}^2\right)^{-1/\kappa}
\)，
移项即得所示上界。
对数形式由 \(E_{\mathrm{int}}=-\log(\det\mathsf P/\prod_j p_j)\) 与 \(\log(p_m/\pi_m)=-\log\pi_m+\log p_m\) 立得。证毕。
\end{proof}


% (User input window)

\begin{remark}[条件通量即条件体积：\texorpdfstring{$V_{\partial}=\prod_m\sqrt{\pi_m}$}{Vpartial=prod sqrt(pi_m)}]\label{rem:xi-pick-poisson-schur-conditional-volume}
在命题 \ref{prop:xi-pick-poisson-schur-flux-ledger} 的设定下，定义边界平行体体积
\[
V_{\partial}:=\sqrt{\det\mathsf P}.
\]
则由 Schur 链式分解 \(\det\mathsf P=\prod_{m=1}^{\kappa}\pi_m\) 得到严格乘积律
\[
\boxed{\;
V_{\partial}=\prod_{m=1}^{\kappa}\sqrt{\pi_m},\qquad
S_{\mathrm{gen}}=-\log\det\mathsf P=\sum_{m=1}^{\kappa}(-\log\pi_m).\;}
\]
因此每一增量 \(-\log\pi_m\) 可解释为“在既有缺陷张成空间上的条件体积损失”；并且 \(0<\pi_m\le p_m\) 给出逐步损失由端点自能上界控制的显式夹逼。
\end{remark}



\begin{proposition}[边界插值裕度解释：\texorpdfstring{$\det\mathsf P$}{det P} 与 \texorpdfstring{$\lambda_{\min}(\mathsf P)$}{lambda_min(P)} 作为可验证余量证书]\label{prop:xi-pick-poisson-interpolation-margin-certificate}
令 \(k_a(w)=(1-\overline a\,w)^{-1}\) 为 Hardy 再生核，并令
\[
\psi_j(w):=\frac{1-|a_j|^2}{1+\overline a_j}\,k_{a_j}(w)
=\frac{1-|a_j|^2}{(1+\overline a_j)(1-\overline a_j w)}\in H^2.
\]
则 \(\mathsf P\) 可写作 \(\{\psi_j\}\) 的 Gram 矩阵：
\[
\boxed{\;
\langle \psi_j,\psi_k\rangle_{H^2}=\mathsf P_{jk}.\;}
\]
因此 \(\det\mathsf P\) 等于 \(\{\psi_j\}\) 在 \(H^2\) 中所张成平行体的平方体积（Gram 行列式定理），其非退化性等价于端点锚定评价子空间的线性独立性。
并且
\[
\lambda_{\min}(\mathsf P)
=\min_{\|c\|_2=1}\left\|\sum_{j=1}^{\kappa}c_j\psi_j\right\|_{H^2}^2
\]
给出该评价子空间的最坏“裕度”。
结合推论 \ref{cor:xi-pick-poisson-gap-lower-bound} 的显式下界，即得一条完全几何化的可验证余量证书：其退化机制被严格归因于贴边（\(p_j\downarrow 0\)）与碰撞（\(\rho_{jk}\downarrow 0\)）。
\leanverified{paper\_xi\_pick\_poisson\_interpolation\_margin\_certificate}
\end{proposition}
\begin{proof}[证明要点]
由 \(\langle k_{a_j},k_{a_k}\rangle_{H^2}=(1-\overline a_j a_k)^{-1}\) 与
\(
\psi_j=\frac{1-|a_j|^2}{1+\overline a_j}k_{a_j}
\)，
可得
\[
\langle \psi_j,\psi_k\rangle_{H^2}
=\frac{(1-|a_j|^2)(1-|a_k|^2)}{(1+\overline a_j)(1+a_k)(1-\overline a_j a_k)}
=\mathsf P_{jk}.
\]
Gram 行列式与最小特征值的变分表述为标准线性代数事实。证毕。
\end{proof}

\begin{definition}[锚定容量（离散视界容量）]\label{def:xi-anchored-capacity-capminus1}
令 \(A=\{a_j\}_{j=1}^{\kappa}\subset\mathbb{D}\) 为有限点状缺陷点集（按重数展开），并令 \(\mathsf P(A)\) 为其 Pick--Poisson Gram 矩阵 \eqref{eq:xi-pick-poisson-gram}。
定义锚定容量（离散视界容量）
\[
\boxed{\;
\mathrm{cap}_{-1}(A):=\bigl(\det\mathsf P(A)\bigr)^{\frac{1}{2\kappa}}
=\exp\!\left(-\frac{1}{2\kappa}S_{\mathrm{gen}}(A)\right)\ \in(0,\infty).\;}
\]
\end{definition}

\begin{proposition}[锚定容量的三性质：置换不变、退化判别与并集次乘性]\label{prop:xi-anchored-capacity-three-properties}
\leanverified{paper\_xi\_anchored\_capacity\_three\_properties}
锚定容量 \(\mathrm{cap}_{-1}\) 具有如下性质：
\begin{enumerate}
  \item \textbf{（置换不变）}对任意重排 \(A\mapsto A'\)（保持多重集不变）有 \(\mathrm{cap}_{-1}(A')=\mathrm{cap}_{-1}(A)\)；
  \item \textbf{（单调退化）}若存在一列配置使得某个端点权重 \(p_j=P_{a_j}(-1)\downarrow 0\)（其余点保持在紧域内），或存在碰撞 \(\rho(a_j,a_k)\downarrow 0\)（其余点固定），则 \(\mathrm{cap}_{-1}(A)\to 0\)；
  \item \textbf{（并集次乘性）}对互不交叠并集 \(A\sqcup B\)（按重数计）有
  \[
  \det\mathsf P(A\cup B)\ \le\ \det\mathsf P(A)\,\det\mathsf P(B),
  \]
  从而
  \[
  \boxed{\;
  \mathrm{cap}_{-1}(A\cup B)
  \le
  \mathrm{cap}_{-1}(A)^{\frac{|A|}{|A|+|B|}}\,
  \mathrm{cap}_{-1}(B)^{\frac{|B|}{|A|+|B|}}.\;}
  \]
\end{enumerate}
因此 \(\mathrm{cap}_{-1}\) 在制度意义上给出“点状缺陷可被外部接口稳定分辨”的单标量尺度；其对数势 \(S_{\mathrm{gen}}=-\log\det\mathsf P\) 为离散容量的可加记账势。
\end{proposition}
\begin{proof}[证明要点]
（1）由行列式对指标重排不变性立得。
\par
（2）由行列式分解 \(\det\mathsf P=(\prod_j p_j)(\prod_{j<k}\rho_{jk}^2)\)（命题 \ref{prop:xi-pick-poisson-det-factorization}）立得。
\par
（3）由并集交叉互能恒等式（命题 \ref{prop:xi-pick-poisson-union-submultiplicativity-cross-energy}）知
\(
\det\mathsf P(A\cup B)=\det\mathsf P(A)\det\mathsf P(B)\cdot \prod_{a\in A,b\in B}\rho(a,b)^2
\le \det\mathsf P(A)\det\mathsf P(B)
\)；
再取 \(1/(2(|A|+|B|))\) 次幂即可。证毕。
\end{proof}

\begin{proposition}[锚定 Hardy 最小范数插值：\texorpdfstring{$\mathsf P\succ 0$}{P>0} 当且仅当锚定插值可行]\label{prop:xi-anchored-hardy-interpolation-pick-equivalence}
沿用命题 \ref{prop:xi-pick-poisson-interpolation-margin-certificate} 的锚定核向量 \(\psi_j\in H^2\)，并定义锚定评价算子
\[
\mathcal{E}_{A,-1}:H^2\to\CC^{\kappa},
\qquad
(\mathcal{E}_{A,-1}f)_j:=\langle f,\psi_j\rangle_{H^2}
=\frac{1-|a_j|^2}{1+\overline a_j}\,f(a_j).
\]
给定数据 \(c\in\CC^{\kappa}\)，考虑锚定插值问题：在约束 \(\mathcal{E}_{A,-1}f=c\) 下最小化 \(\|f\|_{H^2}\)。
则以下结论成立：
\begin{enumerate}
  \item \(\mathsf P(A)\succ 0\) 当且仅当对每个 \(c\) 存在唯一最小范数解 \(f_c\in H^2\)；
  \item 当 \(\mathsf P(A)\succ 0\) 时，最小范数满足闭式
  \[
  \boxed{\;
  \|f_c\|_{H^2}^2=c^\ast\,\mathsf P(A)^{-1}\,c;\;}
  \]
  \item 因而插值算子 \(c\mapsto f_c\) 的稳定性常数被 \(\lambda_{\min}(\mathsf P(A))\) 完全封口：
  \[
  \boxed{\;
  \|f_c\|_{H^2}^2\le \lambda_{\min}(\mathsf P(A))^{-1}\,\|c\|_2^2,}
  \]
  且该上界在极端方向上可达（取 \(c\) 为最小特征向量）。
\end{enumerate}
因此 \(\mathsf P\) 同时是“缺陷几何”的 Gram 矩阵与“锚定最小范数插值”的 Pick 矩阵；\(\det\mathsf P\) 与 \(\lambda_{\min}(\mathsf P)\) 分别刻画体积型与裕度型可验证余量。
\end{proposition}
\begin{proof}[证明要点]
\(\mathsf P(A)\succ 0\) 等价于 \(\{\psi_j\}\) 线性无关，从而算子 \(\mathcal{E}_{A,-1}\) 在 \(\Span\{\psi_j\}\) 上满秩。
最小范数问题在 Hilbert 空间中满足 representer 定理：解必形如 \(f_c=\sum_j \alpha_j\psi_j\)。
代入约束得到线性方程 \(\mathsf P(A)\alpha=c\)，故唯一解为 \(\alpha=\mathsf P(A)^{-1}c\)。
于是
\(
\|f_c\|^2=\alpha^\ast \mathsf P(A)\alpha=c^\ast\mathsf P(A)^{-1}c
\)。
稳定性界由 Rayleigh 商估计立得。证毕。
\end{proof}

\begin{proposition}[最小证书阶数的插值解释：Toeplitz 显影门槛等价于锚定插值条件数门槛]\label{prop:xi-toeplitz-threshold-via-anchored-interpolation}
设 \(A=\{a_j\}_{j=1}^{\kappa}\subset\mathbb{D}\) 处于分离域内并令 \(h_{\min}:=\min_j(1-|a_j|)\)。
令 \(f_c\) 为命题 \ref{prop:xi-anchored-hardy-interpolation-pick-equivalence} 的最小范数解，并令 \(\Pi_N:H^2\to\Span\{1,w,\dots,w^N\}\) 为次数截断正交投影，记 \(f_c^{(N)}:=\Pi_N f_c\)。
则存在常数 \(C>0\)（依赖至多于 \(\kappa\)）使得：若
\[
\boxed{\;
N\ \ge\ C\,\frac{\log\kappa(\mathsf P(A))}{h_{\min}},\qquad
\kappa(\mathsf P):=\frac{\lambda_{\max}(\mathsf P)}{\lambda_{\min}(\mathsf P)},\;}
\]
则对所有 \(c\in\CC^{\kappa}\) 有
\[
\boxed{\;
\|f_c-f_c^{(N)}\|_{H^2}\ \le\ \frac12\,\|f_c\|_{H^2}.\;}
\]
从而在该门槛之后，有限阶 Toeplitz 证书在缺陷核张成子空间上的负谱裕度与 \(\mathsf P\) 的谱裕度同阶并稳定显影（与命题 \ref{prop:xi-toeplitz-negative-spectrum-margin-threshold} 与推论 \ref{cor:xi-toeplitz-visibility-complexity-threshold} 的门槛律一致）。
因此“证书阶数不足导致缺陷不可见”在严格意义上等价于“锚定插值条件数过大导致截断不可控”。
\end{proposition}
\begin{proof}[证明要点]
由 \(k_{a}(w)=\sum_{n\ge 0}\overline a^{\,n}w^n\) 得 \(\|\Pi_{>N}k_a\|_{H^2}=|a|^{N+1}(1-|a|^2)^{-1/2}\)。
对锚定向量 \(\psi_j=\frac{1-|a_j|^2}{1+\overline a_j}k_{a_j}\) 得
\(
\|\Pi_{>N}\psi_j\|_{H^2}\le \sqrt{p_j}\,|a_j|^{N+1}\le \sqrt{p_j}\,e^{-Nh_{\min}}
\)。
又 \(f_c=\sum_j\alpha_j\psi_j\) 且 \(\alpha=\mathsf P^{-1}c\)，从而 \(\|\alpha\|_2\le \lambda_{\min}(\mathsf P)^{-1}\|c\|_2\)。
将尾部估计与 \(\kappa(\mathsf P)=\lambda_{\max}/\lambda_{\min}\) 合并即可得到所示“对所有 \(c\)”的相对误差阈值；常数 \(C\) 吸收 \(\kappa\) 维度因子与 \(p_\Sigma=\mathrm{tr}\,\mathsf P\) 的上界（推论 \ref{cor:xi-pick-poisson-condition-number-certificate}）。证毕。
\end{proof}

\begin{proposition}[点状缺陷的确定性 \texorpdfstring{$\kappa$}{k}-DPP：以 \texorpdfstring{$\det\mathsf P$}{det P} 为密度的最小排斥分布]\label{prop:xi-pick-poisson-kdpp}
取一有限候选网格 \(\Omega=\{x_1,\dots,x_M\}\subset\mathbb{D}\)（或其上半平面像）。
定义 Pick--Poisson 正定核
\[
K_{-1}(w,z):=\frac{(1-|w|^2)(1-|z|^2)}{(1+\overline w)(1+z)(1-\overline w z)}\qquad(w,z\in\mathbb{D}),
\]
并令 \(L\in\Mat_M(\CC)_{\mathrm{Herm}}\) 为其在 \(\Omega\) 上的 Gram：
\[
L_{uv}:=K_{-1}(x_u,x_v)\qquad(1\le u,v\le M).
\]
则 \(L\succeq 0\)，且对任意 \(S\subseteq\Omega\) 有 \(\det(L_S)=\det\mathsf P(S)\)。
对固定点数 \(\kappa\) 与参数 \(\beta>0\)，在 \(\Omega\) 的 \(\kappa\)-子集上定义分布
\[
\mathbb{P}_{\beta}^{(\Omega)}(S)\ \propto\ \det(L_S)^{\beta}\qquad(|S|=\kappa).
\]
则当 \(\beta=1\) 时，上述分布即为标准 \(\kappa\)-DPP（fixed-size determinantal distribution；亦即 \(k\)-DPP/\(L\)-ensemble），并可看作 \(L\)-ensemble DPP 在条件于点数为 \(\kappa\) 后得到的 fixed-size 版本（参见 \cite{KuleszaTaskar2012DPP}）。
并且其对数密度等于 \(-S_{\mathrm{gen}}(S)\)（差一个只依赖于 \(\Omega\) 的归一化常数），从而给出“最大化可验证余量 \(\Longleftrightarrow\) 最小化镜像能量”的确定性排斥实现。
\end{proposition}
\begin{proof}[证明要点]
对有限 \(\Omega\)，矩阵 \(L\succeq 0\) 定义了一个 \(k\)-DPP，其在大小为 \(\kappa\) 的子集上的权重正比于 \(\det(L_S)\)（\cite{KuleszaTaskar2012DPP}）。
而 \(S_{\mathrm{gen}}(S)=-\log\det\mathsf P(S)=-\log\det(L_S)\)（定义 \ref{def:xi-generalized-defect-entropy-logdet}），故得到对数密度表述。证毕。
\end{proof}

\begin{proposition}[均场极限模板：\texorpdfstring{$\kappa^{-2}S_{\mathrm{gen}}$}{k^{-2}S_gen} 的主阶由镜像 Green 能量泛函决定]\label{prop:xi-pick-poisson-meanfield-template}
在命题 \ref{prop:xi-pick-poisson-image-charge-green-energy} 的上半平面坐标下，令 \(t_1,\dots,t_\kappa\in\mathcal{K}\subset\mathbb{H}\)（\(\mathcal{K}\) 紧致且 \(\Im t\ge y_0>0\)）。
设经验测度 \(\mu_\kappa:=\frac{1}{\kappa}\sum_{j=1}^{\kappa}\delta_{t_j}\) 弱收敛到 \(\mu\in\mathcal{P}(\mathcal{K})\)，且点配置在尺度上不发生宏观塌缩（例如存在统一分离下界 \(\rho_{\min}\ge\rho_0>0\)）。
则有主阶极限
\[
\boxed{\;
\lim_{\kappa\to\infty}\frac{1}{\kappa^2}S_{\mathrm{gen}}(t_1,\dots,t_\kappa)
=\iint_{\mathcal{K}\times\mathcal{K}}G_{\mathbb{H}}(t,s)\,\dd\mu(t)\dd\mu(s),\;}
\]
其中 \(G_{\mathbb{H}}(t,s)=\log\left|\frac{t-\overline s}{t-s}\right|\)。
自能项 \(\kappa^{-2}\sum_j(-\log\Im t_j)\) 在固定 \(\mathcal{K}\) 内为次阶并退化为 \(0\)。
\end{proposition}
\begin{proof}[证明要点]
在分离域内 \(G_{\mathbb{H}}\) 在 \(\mathcal{K}\times\mathcal{K}\) 上有界连续，故
\(
\frac{2}{\kappa^2}\sum_{j<k}G_{\mathbb{H}}(t_j,t_k)
\to \iint G_{\mathbb{H}}\,\dd\mu\dd\mu
\)；
而自能项为 \(O(\kappa)\) 故除以 \(\kappa^2\) 消失。证毕。
\end{proof}

\begin{corollary}[平衡测度的 Euler--Lagrange 方程：视界锚定导致的边界常势条件]\label{cor:xi-pick-poisson-equilibrium-euler-lagrange}
在命题 \ref{prop:xi-pick-poisson-meanfield-template} 的连续能量泛函
\(
\mathcal{E}(\mu):=\iint G_{\mathbb{H}}(t,s)\,\dd\mu(t)\dd\mu(s)
\)
上，若 \(\mu^\ast\) 为 \(\mathcal{P}(\mathcal{K})\) 上的极小平衡测度，则存在常数 \(C\) 使其势函数
\[
U^{\mu^\ast}(t):=\int_{\mathcal{K}}G_{\mathbb{H}}(t,s)\,\dd\mu^\ast(s)
\]
满足 Euler--Lagrange 条件
\[
\boxed{\;
U^{\mu^\ast}(t)=C\ \text{于}\ \mathrm{supp}(\mu^\ast),\qquad
U^{\mu^\ast}(t)\ge C\ \text{于}\ \mathcal{K}.\;}
\]
该常势条件给出连续极限中“最可见缺陷分布”的变分学刻画（参见对数势理论中的 Frostman 原理，\cite{SaffTotik1997LogarithmicPotentials}）。
\end{corollary}
\begin{proof}[证明要点]
这是对数势能量在紧集上的标准平衡测度结论：极小性等价于势函数满足常势/障碍条件。证毕。
\end{proof}

\begin{proposition}[负谱体积的集中不等式：分离域内 \texorpdfstring{$\log\det\mathsf P$}{log det P} 的 Lipschitz 性与 McDiarmid 型集中]\label{prop:xi-pick-poisson-logdet-concentration}
设点配置满足统一分离下界 \(\rho_{\min}\ge\rho_0>0\) 与高度下界 \(\Im t_j\ge y_0>0\)，并令
\(
F(t_1,\dots,t_\kappa):=\log\det\mathsf P(t_1,\dots,t_\kappa)
\)。
则 \(F\) 在该分离域上为全局 Lipschitz（常数仅依赖 \(\kappa,\rho_0,y_0\)）。
因此对独立采样 \(t_1,\dots,t_\kappa\)（或满足有界差分条件的马氏链采样）具有 McDiarmid 型集中：存在常数 \(c>0\) 使对任意 \(u>0\),
\[
\boxed{\;
\mathbb{P}\bigl(|F-\mathbb{E}F|\ge u\bigr)\ \le\ 2\exp\!\left(-c\,\frac{u^2}{\kappa}\right)\;}
\]
（参见 \cite{McDiarmid1989BoundedDifferences}）。
\end{proposition}
\begin{proof}[证明要点]
由 \(F=-S_{\mathrm{gen}}\) 且命题 \ref{prop:xi-pick-poisson-energy-balance-equations} 给出 \(\nabla S_{\mathrm{gen}}\) 的显式分式形式；在分离域内分母一致远离零，故梯度一致有界，从而 Lipschitz。
有界差分/独立采样下的集中不等式由 McDiarmid 方法推出。证毕。
\end{proof}

\begin{proposition}[视界粗粒化下的自由能耗散：Poisson 半群对相对熵的严格收缩]\label{prop:xi-pick-poisson-poisson-semigroup-free-energy-dissipation}
令 \(T_t=\exp(-t|D|)\) 为 \(\RR\) 上 Poisson 半群（命题 \ref{con:discussion-dtn-generator}），并令其在 \(\RR^\kappa\) 上的张量化半群为 \(T_t^{\otimes \kappa}\)。
对任意两概率测度 \(\mu,\nu\)（定义在 \(\RR^\kappa\) 上），记
\[
\mu_t:=\mu\,T_t^{\otimes\kappa},\qquad \nu_t:=\nu\,T_t^{\otimes\kappa},
\]
即分别在每个坐标上与 Poisson 核卷积的 Markov 推前。
则相对熵满足半群单调性
\[
\boxed{\;
D(\mu_t\|\nu_t)\ \le\ D(\mu_s\|\nu_s)\qquad(t\ge s\ge 0).\;}
\]
若进一步选取某个对 \(T_t^{\otimes\kappa}\) 不变的参考测度 \(\nu\)（例如在周期化/有限体积口径下取 Haar/均匀基准），则 \(t\mapsto D(\mu_t\|\nu)\) 单调不增；
并且在平滑条件下其导数由生成元 \(\sum_{j=1}^{\kappa}|D_{\gamma_j}|\) 的 Dirichlet 形式给出：
\[
\frac{\dd}{\dd t}D(\mu_t\|\nu)=-\mathcal{I}(\mu_t)\le 0,
\]
其中 \(\mathcal{I}\) 为对应的非局域 Fisher 信息。
因此“视界回收/粗粒化”在多缺陷统计系综层面仍表现为严格的热力学势耗散：粗粒化并非近似操作，而由几何生成元唯一确定的不可逆流。
\end{proposition}
\begin{proof}[证明要点]
由半群性质 \(\mu_t=(\mu_s)T_{t-s}^{\otimes\kappa}\)、\(\nu_t=(\nu_s)T_{t-s}^{\otimes\kappa}\)，
以及相对熵对 Markov 信道的收缩性（数据处理不等式）得
\(
D(\mu_t\|\nu_t)\le D(\mu_s\|\nu_s)
\)。
当 \(\nu\) 为不变测度时，\(\nu_t=\nu\) 并得到 \(D(\mu_t\|\nu)\) 的单调性；导数公式为标准熵耗散恒等式（以生成元表示）。证毕。
\end{proof}

\begin{conjecture}[两源接口的能量--指数一致性：去链接缺口与交叉互能的量级同构]\label{conj:xi-energy-index-gap-comparability}
设两源接口的 Jones 交换子标量为 \(\lambda=\exp(-\Gamma)\)（命题 \ref{prop:op-algebra-jones-scalar-equals-exp-minus-gap}），并记指数缺口为 \(\Gamma\ge 0\)。
对几何并集 \(A\sqcup B\) 定义交叉互能
\[
\mathcal{E}_{\times}(A,B):=2\sum_{a\in A}\sum_{b\in B}\bigl(-\log\rho(a,b)\bigr)
=-\log\frac{\det\mathsf P(A\cup B)}{\det\mathsf P(A)\det\mathsf P(B)}.
\]
在稳定分离域内（\(\rho(a,b)\) 远离 \(0\) 且端点通量不塌缩），存在常数 \(c_1,c_2>0\) 使得
\[
\boxed{\;
c_1\,\Gamma(A,B)\ \le\ \mathcal{E}_{\times}(A,B)\ \le\ c_2\,\Gamma(A,B),\;}
\]
其中 \(\Gamma(A,B)\) 表示由相应 Jones 投影交换子标量定义的指数缺口。
该对齐猜想表明：在稳定域内，“去链接所需资源”（指数缺口）与“并集互能”（负谱交叉项）在数量级上同构，且二者均由同一分离参数控制。
\end{conjecture}

\begin{proposition}[锚定容量的算子范数表述：\texorpdfstring{$\mathrm{cap}_{-1}$}{cap-1} 为插值右逆外乘范数的几何均值]\label{prop:xi-anchored-capacity-exterior-power-interpolant}
\leanverified{paper\_xi\_anchored\_capacity\_exterior\_power\_interpolant}
沿用命题 \ref{prop:xi-anchored-hardy-interpolation-pick-equivalence} 的锚定评价算子
\(\mathcal{E}_{A,-1}:H^2\to\CC^{\kappa}\) 与 Gram 矩阵 \(\mathsf P(A)\succ0\)。
记由最小范数原则定义的\emph{锚定插值右逆}（最小范数插值算子）
\[
\mathcal{I}_{A,-1}:\CC^{\kappa}\to H^2,\qquad \mathcal{I}_{A,-1}(c):=f_c,
\]
其中 \(f_c\) 为满足 \(\mathcal{E}_{A,-1}f=c\) 的唯一最小范数解。
则 \(\mathcal{I}_{A,-1}\) 线性有界并满足
\[
\boxed{\;
\mathcal{I}_{A,-1}^{\ast}\mathcal{I}_{A,-1}=\mathsf P(A)^{-1},
\qquad
\det(\mathcal{I}_{A,-1}^{\ast}\mathcal{I}_{A,-1})=\det\mathsf P(A)^{-1}.\;}
\]
从而锚定容量
\(\mathrm{cap}_{-1}(A):=\bigl(\det\mathsf P(A)\bigr)^{\frac{1}{2\kappa}}\)
可被完全转写为一个标准算子不变量：
\[
\boxed{\;
\mathrm{cap}_{-1}(A)^{-1}
=\bigl\|\wedge^{\kappa}\mathcal{I}_{A,-1}\bigr\|^{1/\kappa}\;}
\]
其中 \(\wedge^{\kappa}\mathcal{I}_{A,-1}\) 表示诱导于 \(\kappa\) 次外幂的算子，而其范数等于 \(\mathcal{I}_{A,-1}\) 的 \(\kappa\) 个奇异值的乘积。
\end{proposition}
\begin{proof}[证明要点]
由命题 \ref{prop:xi-anchored-hardy-interpolation-pick-equivalence}，
对任意 \(c\) 有 \(\|f_c\|_{H^2}^2=c^\ast\mathsf P(A)^{-1}c\)。
将 \(f_c=\mathcal{I}_{A,-1}c\) 代入即得
\(\langle c,\mathcal{I}_{A,-1}^{\ast}\mathcal{I}_{A,-1}c\rangle
=c^\ast\mathsf P(A)^{-1}c\)，从而 \(\mathcal{I}_{A,-1}^{\ast}\mathcal{I}_{A,-1}=\mathsf P(A)^{-1}\)。
外幂范数恒等式 \(\|\wedge^{\kappa}\mathcal{I}\|^2=\det(\mathcal{I}^{\ast}\mathcal{I})\) 为有限秩算子的标准奇异值事实，代入即得所示结论。证毕。
\end{proof}

\begin{corollary}[视界稳定子群的共变律：\texorpdfstring{$\mathrm{cap}_{-1}$}{cap-1} 的尺度漂移由 \texorpdfstring{$\phi'(-1)$}{phi'(-1)} 封口]\label{cor:xi-anchored-capacity-stabilizer-covariance}
\leanverified{paper\_xi\_anchored\_capacity\_stabilizer\_covariance}
令 \(\phi\) 为单位圆盘自同构且 \(\phi(-1)=-1\)。
记 \(\phi'(-1)>0\) 为其在 \(-1\) 处的 Julia--Carath\'eodory 角导数。
则伪双曲距离不变：\(\rho(\phi(a),\phi(b))=\rho(a,b)\)。
并且端点 Poisson 权重满足
\[
\boxed{\;P_{\phi(a)}(-1)=\phi'(-1)\,P_a(-1)\;}
\]
从而对任意有限点集 \(A=\{a_j\}_{j=1}^{\kappa}\subset\mathbb{D}\) 有协变律
\[
\boxed{\;
\det\mathsf P(\phi(A))=\phi'(-1)^{\kappa}\det\mathsf P(A),\qquad
\mathrm{cap}_{-1}(\phi(A))=\phi'(-1)^{1/2}\,\mathrm{cap}_{-1}(A).\;}
\]
因此 \(-\log\det\mathsf P\) 的唯一同调漂移为 \(-\kappa\log\phi'(-1)\)，而所有互能项（由 \(\rho\) 给出）严格不变。
\end{corollary}
\begin{proof}[证明要点]
该结论是命题 \ref{prop:xi-pick-poisson-det-scaling-law} 的 disk--stabilizer 形式：在 \(\phi(-1)=-1\) 的稳定子群上，
\(\rho\) 不变而 \(P_a(-1)\) 以 \(\phi'(-1)\) 协变；再用行列式分解
\(\det\mathsf P=(\prod_j P_{a_j}(-1))(\prod_{j<k}\rho(a_j,a_k)^2)\)（命题 \ref{prop:xi-pick-poisson-det-factorization}）即可。证毕。
\end{proof}

\begin{theorem}[锚定插值右逆的 Schatten--容量下界与各向同性判别]\label{thm:xi-anchored-schatten-capacity-rigidity}
\leanverified{paper\_xi\_anchored\_schatten\_capacity\_rigidity}
沿用命题 \ref{prop:xi-anchored-capacity-exterior-power-interpolant} 的记号，并令
\[
s_1,\dots,s_\kappa>0
\]
为锚定插值右逆 \(\mathcal{I}_{A,-1}\) 的奇异值。
则对任意 \(p\in(0,\infty]\) 都有
\[
\boxed{\;
\|\mathcal{I}_{A,-1}\|_{S_p}
\ge
\kappa^{1/p}\,\mathrm{cap}_{-1}(A)^{-1},\;}
\]
其中 \(p=\infty\) 时约定 \(\kappa^{1/\infty}=1\) 且 \(\|\cdot\|_{S_\infty}=\|\cdot\|_{\mathrm{op}}\)。
并且等号成立当且仅当
\[
s_1=\cdots=s_\kappa=\mathrm{cap}_{-1}(A)^{-1},
\]
等价地，
\[
\boxed{\;\mathsf P(A)=\mathrm{cap}_{-1}(A)^2\,I_\kappa.\;}
\]
\end{theorem}
\begin{proof}
由命题 \ref{prop:xi-anchored-capacity-exterior-power-interpolant}，
\[
\prod_{j=1}^{\kappa}s_j
=
\bigl\|\wedge^\kappa \mathcal I_{A,-1}\bigr\|
=
\mathrm{cap}_{-1}(A)^{-\kappa},
\]
故奇异值的几何均值为
\[
\Bigl(\prod_{j=1}^{\kappa}s_j\Bigr)^{1/\kappa}
=
\mathrm{cap}_{-1}(A)^{-1}.
\]
若 \(0<p<\infty\)，对非负数列 \((s_j^p)_{j=1}^{\kappa}\) 应用 AM--GM，不等式给出
\[
\frac1\kappa\sum_{j=1}^{\kappa}s_j^p
\ge
\Bigl(\prod_{j=1}^{\kappa}s_j^p\Bigr)^{1/\kappa}
=
\mathrm{cap}_{-1}(A)^{-p},
\]
从而
\[
\|\mathcal I_{A,-1}\|_{S_p}^p
=
\sum_{j=1}^{\kappa}s_j^p
\ge
\kappa\,\mathrm{cap}_{-1}(A)^{-p}.
\]
若 \(p=\infty\)，则
\[
\|\mathcal I_{A,-1}\|_{\mathrm{op}}
=
\max_{1\le j\le \kappa}s_j
\ge
\Bigl(\prod_{j=1}^{\kappa}s_j\Bigr)^{1/\kappa}
=
\mathrm{cap}_{-1}(A)^{-1}.
\]
两种情形下的等号条件都等价于全部奇异值相等。
再由
\(
\mathcal I_{A,-1}^\ast\mathcal I_{A,-1}=\mathsf P(A)^{-1}
\)
可知，奇异值全相等当且仅当
\[
\mathsf P(A)^{-1}=\mathrm{cap}_{-1}(A)^{-2}\,I_\kappa,
\]
亦即
\(
\mathsf P(A)=\mathrm{cap}_{-1}(A)^2\,I_\kappa
\)。
证毕。
\end{proof}

\begin{corollary}[锚定容量的极值原理：固定端点通量下的最优上界与必要结构]\label{cor:xi-anchored-capacity-extremal-fixed-flux}
\leanverified{paper\_xi\_anchored\_capacity\_extremal\_fixed\_flux}
固定 \(\kappa\) 并给定总端点通量 \(p_{\Sigma}:=\sum_{j=1}^{\kappa}p_j\)（其中 \(p_j=P_{a_j}(-1)\)）。
则
\[
\boxed{\;
\det\mathsf P(A)\ \le\ \Bigl(\frac{p_{\Sigma}}{\kappa}\Bigr)^{\kappa},\qquad
\mathrm{cap}_{-1}(A)\ \le\ \Bigl(\frac{p_{\Sigma}}{\kappa}\Bigr)^{1/2}.\;}
\]
并且若存在一列构型使得 \(\det\mathsf P(A)\) 逼近上述上界，则必有
（i）权重渐近均匀：\(p_j\to p_{\Sigma}/\kappa\)（按重数）；
（ii）点对渐近极限分离：\(\rho(a_j,a_k)\to 1\)（所有 \(j\neq k\)）。
换言之，“最可验证”构型在制度允许的几何自由度中由\emph{自能均匀化}与\emph{互能分离化}共同强制。
\end{corollary}
\begin{proof}[证明要点]
由 \(\rho\le 1\) 与行列式分解得 \(\det\mathsf P\le\prod_j p_j\)；
再用 AM--GM 得 \(\prod_j p_j\le (p_{\Sigma}/\kappa)^{\kappa}\)，从而得到上界（推论 \ref{cor:xi-pick-poisson-energy-fixed-flux-extremal} 的同型表述）。
若逼近上界，则必须同时逼近 \(\rho\)-项的上界与 AM--GM 的等号条件，给出（i）与（ii）。证毕。
\end{proof}

\begin{theorem}[固定总端点通量下行列式前沿的非达成性]\label{thm:xi-anchored-fixed-flux-frontier-nonattainment}
\leanverified{paper\_xi\_anchored\_fixed\_flux\_frontier\_nonattainment}
固定 \(\kappa\ge 2\) 与总端点通量
\[
p_{\Sigma}:=\sum_{j=1}^{\kappa}p_j,
\qquad
p_j:=P_{a_j}(-1).
\]
则有
\[
\boxed{\;
\sup_{\substack{|A|=\kappa\\ \sum_j p_j=p_{\Sigma}}}\det\mathsf P(A)
=
\Bigl(\frac{p_{\Sigma}}{\kappa}\Bigr)^\kappa,\;}
\]
但该上确界不被任何有限构型 \(A\subset\mathbb D\) 达到。
更精确地，若 \(\{A^{(\nu)}\}_{\nu\ge 1}\) 为任意逼近上确界的构型序列，则必有
\[
p_j^{(\nu)}-\frac{p_{\Sigma}}{\kappa}\longrightarrow 0
\qquad(1\le j\le \kappa),
\]
并且
\[
\rho\bigl(a_j^{(\nu)},a_k^{(\nu)}\bigr)\longrightarrow 1
\qquad(j\neq k).
\]
\end{theorem}
\begin{proof}
上界由推论 \ref{cor:xi-anchored-capacity-extremal-fixed-flux} 直接给出。
另一方面，对任意有限构型，
\[
\det\mathsf P(A)
=
\Bigl(\prod_{j=1}^{\kappa}p_j\Bigr)\prod_{1\le j<k\le \kappa}\rho(a_j,a_k)^2.
\]
当 \(\kappa\ge 2\) 且点互异时，恒有 \(0<\rho(a_j,a_k)<1\)（\(j\neq k\)），故
\[
\det\mathsf P(A)
<
\prod_{j=1}^{\kappa}p_j
\le
\Bigl(\frac{p_{\Sigma}}{\kappa}\Bigr)^\kappa.
\]
于是上确界不能在任何有限构型上达到。

为说明该上界可逼近，取命题 \ref{prop:xi-pick-poisson-image-charge-green-energy} 的上半平面规一化坐标
\[
t_j^{(\nu)}:=x_j^{(\nu)}+\mathrm{i}\,\frac{p_{\Sigma}}{\kappa},
\]
其中 \(|x_j^{(\nu)}-x_k^{(\nu)}|\to\infty\)（\(j\neq k\)）。
记 \(a_j^{(\nu)}\in\mathbb D\) 为对应圆盘点。
则
\(
P_{a_j^{(\nu)}}(-1)=\Im(t_j^{(\nu)})=p_{\Sigma}/\kappa
\)
且
\(
\rho(a_j^{(\nu)},a_k^{(\nu)})\to 1
\)，
从而
\[
\det\mathsf P(A^{(\nu)})
\longrightarrow
\Bigl(\frac{p_{\Sigma}}{\kappa}\Bigr)^\kappa.
\]

最后，若 \(\det\mathsf P(A^{(\nu)})\to (p_{\Sigma}/\kappa)^\kappa\)，则
\[
1\ge
\frac{\prod_j p_j^{(\nu)}}{(p_{\Sigma}/\kappa)^\kappa}
\ge
\frac{\det\mathsf P(A^{(\nu)})}{(p_{\Sigma}/\kappa)^\kappa}
\longrightarrow 1,
\]
故
\(
\prod_j p_j^{(\nu)}/(p_{\Sigma}/\kappa)^\kappa\to 1
\)。
由 AM--GM 的等号稳定性，得到
\(
p_j^{(\nu)}-p_{\Sigma}/\kappa\to 0
\)
对每个 \(j\) 成立。
再由
\[
\prod_{1\le j<k\le \kappa}\rho(a_j^{(\nu)},a_k^{(\nu)})^2
=
\frac{\det\mathsf P(A^{(\nu)})}{\prod_j p_j^{(\nu)}}
\longrightarrow 1
\]
且各因子均不超过 \(1\)，可知每一对 \(j\neq k\) 都满足
\(
\rho(a_j^{(\nu)},a_k^{(\nu)})\to 1
\)。
证毕。
\end{proof}

\begin{corollary}[锚定容量极值的定量稳定性：权重均匀化与最小分离度的显式误差界（51）]\label{cor:xi-anchored-capacity-extremal-stability-quant}
\leanverified{paper\_xi\_anchored\_capacity\_extremal\_stability\_quant}
令 \(A=\{a_1,\dots,a_\kappa\}\subset\mathbb{D}\) 为锚定构型，并记端点通量 \(p_j:=P_{a_j}(-1)\)、\(p_{\Sigma}:=\sum_{j=1}^{\kappa}p_j\)。\footnote{据内部文稿 \cite{Internal2025ResolutionFoldingPhiPi}（2025）整理。}
定义极值亏损
\[
\boxed{\;
\varepsilon(A):=\kappa\log\!\Bigl(\frac{p_{\Sigma}}{\kappa}\Bigr)-\log\det\mathsf P(A)\ \ge\ 0\;}
\]
则成立如下定量稳定估计：
\begin{enumerate}
  \item 令 \(q_j:=p_j/p_{\Sigma}\)（\(\sum_j q_j=1\)）与均匀分布 \(u_j:=1/\kappa\)。则有
  \[
  \boxed{\;
  \sum_{j=1}^{\kappa}|q_j-u_j|
  \ \le\ 
  \sqrt{\frac{2\,\varepsilon(A)}{\kappa}},\qquad
  \sum_{j=1}^{\kappa}\Bigl|p_j-\frac{p_{\Sigma}}{\kappa}\Bigr|
  \ \le\
  p_{\Sigma}\sqrt{\frac{2\,\varepsilon(A)}{\kappa}}\;}
  \]
  从而小亏损强迫端点权重在总变差意义下逼近均匀化。
  \item 记最小分离度 \(\rho_{\min}(A):=\min_{j\ne k}\rho(a_j,a_k)\in(0,1]\)。则有
  \[
  \boxed{\;
  \rho_{\min}(A)\ \ge\ \exp\!\Bigl(-\frac{\varepsilon(A)}{\kappa(\kappa-1)}\Bigr).\;}
  \]
  特别地，若 \(\varepsilon(A)=o(\kappa^2)\)，则必有 \(\rho_{\min}(A)\to 1\)，即极值逼近强迫点对渐近分离。
\end{enumerate}
\end{corollary}

\begin{proof}
由行列式分解（命题 \ref{prop:xi-pick-poisson-det-factorization}）
\[
\log\det\mathsf P(A)
=\sum_{j=1}^{\kappa}\log p_j+2\sum_{1\le j<k\le \kappa}\log\rho(a_j,a_k),
 \qquad \rho(a_j,a_k)\in(0,1],
\]
从而
\[
\varepsilon(A)
=\underbrace{\kappa\log\!\Bigl(\frac{p_{\Sigma}}{\kappa}\Bigr)-\sum_{j=1}^{\kappa}\log p_j}_{=:\,\varepsilon_w\ \ge\ 0}
\;+\;
\underbrace{\Bigl(-2\sum_{j<k}\log\rho(a_j,a_k)\Bigr)}_{=:\,\varepsilon_\rho\ \ge\ 0},
\qquad \varepsilon_w,\varepsilon_\rho\le \varepsilon(A).
\]
令 \(q_j=p_j/p_{\Sigma}\)、\(u_j=1/\kappa\)。则
\[
\varepsilon_w=\sum_{j=1}^{\kappa}\log\!\Bigl(\frac{p_{\Sigma}/\kappa}{p_j}\Bigr)
=\sum_{j=1}^{\kappa}\log\!\Bigl(\frac{u_j}{q_j}\Bigr)
=\kappa\,D(u\Vert q),
\]
其中 \(D(u\Vert q)=\sum_j u_j\log(u_j/q_j)\) 为 KL 散度。由 Pinsker 不等式
\(\|u-q\|_1\le\sqrt{2D(u\Vert q)}\)
得
\[
\sum_{j=1}^{\kappa}|q_j-u_j|
\le
\sqrt{\frac{2\,\varepsilon_w}{\kappa}}
\le
\sqrt{\frac{2\,\varepsilon(A)}{\kappa}},
\]
乘以 \(p_{\Sigma}\) 即得第一条的第二个不等式。
\par
对第二条，由 \(\rho(a_j,a_k)\ge\rho_{\min}(A)\) 得
\[
\varepsilon_\rho
=-2\sum_{j<k}\log\rho(a_j,a_k)
\ge
-2\binom{\kappa}{2}\log\rho_{\min}(A)
=\kappa(\kappa-1)\bigl(-\log\rho_{\min}(A)\bigr).
\]
又 \(\varepsilon_\rho\le \varepsilon(A)\)，故
\(-\log\rho_{\min}(A)\le \varepsilon(A)/(\kappa(\kappa-1))\)，即得到所示下界。证毕。
\end{proof}

\begin{proposition}[交叉互能的高斯互信息表述：并集缺口等于对数行列式互信息]\label{prop:xi-cross-energy-gaussian-mutual-information}
\leanverified{paper\_xi\_cross\_energy\_gaussian\_mutual\_information}
令 \(A\sqcup B\) 为互不交叠并集（按重数计），并将 \(\mathsf P(A\cup B)\) 视为分块协方差矩阵。
取\emph{适当的}复高斯向量 \(Z_{A\cup B}\sim\mathcal N_{\CC}(0,\mathsf P(A\cup B))\)，并令 \(Z_A,Z_B\) 为其在两块上的分量。
定义并集缺口（交叉互能）
\[
\mathcal I(A;B):=-\log\det\mathsf P(A\cup B)+\log\det\mathsf P(A)+\log\det\mathsf P(B)\ \ge\ 0.
\]
则有恒等式
\[
\boxed{\;
\mathcal I(A;B)=I(Z_A;Z_B),\;}
\]
其中右端为（复）高斯互信息（以自然对数计）。
因此交叉互能不仅是几何量（命题 \ref{prop:xi-pick-poisson-union-submultiplicativity-cross-energy}），亦是“二源残余可链接性”的信息论精确度量。
\end{proposition}
\begin{proof}[证明要点]
对适当复高斯向量有微分熵公式 \(h(Z)=\log\det(\pi e\,\Sigma)\)（\(\Sigma\) 为协方差）。
故
\(
I(Z_A;Z_B)=h(Z_A)+h(Z_B)-h(Z_{A\cup B})
=\log\det\mathsf P(A)+\log\det\mathsf P(B)-\log\det\mathsf P(A\cup B)
\)，即得所示恒等式。证毕。
\end{proof}

\begin{proposition}[\texorpdfstring{$\beta$}{beta}-视界系综的主阶自由能：\texorpdfstring{$\kappa^2$}{k^2} 级别由镜像 Green 能量决定]\label{prop:xi-beta-ensemble-leading-free-energy}
在紧致可行域 \(\mathcal K\subset\mathbb H\) 上考虑 \(\beta\)-系综
\[
\mathbb P_{\kappa,\beta}(\dd t_1\cdots \dd t_\kappa)
\ \propto\
\exp\!\bigl(-\beta\,S_{\mathrm{gen}}(t_1,\dots,t_\kappa)\bigr)\,\dd t_1\cdots \dd t_\kappa,
\]
其中 \(S_{\mathrm{gen}}\) 为命题 \ref{prop:xi-pick-poisson-image-charge-green-energy} 给出的镜像 Green 能量（含自能项）。
记配分函数为 \(Z_{\kappa,\beta}\)。
则在“不发生宏观塌缩”的稳定口径下，其主阶自由能满足形式极限
\[
\boxed{\;
\lim_{\kappa\to\infty}\frac{1}{\kappa^2}\log Z_{\kappa,\beta}
=-\beta\inf_{\mu\in\mathcal P(\mathcal K)}
\iint_{\mathcal K\times\mathcal K}G_{\mathbb H}(t,s)\,\dd\mu(t)\dd\mu(s)\;}
\]
其中 \(G_{\mathbb H}(t,s)=\log\left|\frac{t-\overline s}{t-s}\right|\)。
并且极小平衡测度由常势 Euler--Lagrange 条件唯一刻画（推论 \ref{cor:xi-pick-poisson-equilibrium-euler-lagrange}）。
\end{proposition}
\begin{proof}[证明要点]
由命题 \ref{prop:xi-pick-poisson-meanfield-template}，\(\kappa^{-2}S_{\mathrm{gen}}\) 的主阶由双重积分能量给出，而自能项仅为 \(O(\kappa)\)。
由于 \(\log\mathrm{Vol}(\mathcal K^\kappa)=O(\kappa)\)，体积熵贡献在 \(\kappa^{-2}\) 归一化下消失；
从而主阶由能量的极小值支配（Varadhan/拉普拉斯原理的主阶形式）。证毕。
\end{proof}

\begin{conjecture}[线性统计的中心极限定律：锚定确定性点过程的涨落由 \texorpdfstring{$H^{1/2}$}{H^{1/2}} 半范数控制]\label{conj:xi-kdpp-linear-statistics-clt}
在命题 \ref{prop:xi-pick-poisson-kdpp} 的 \(\kappa\)-DPP 口径下，令 \(\kappa\to\infty\) 并在分离域与高度下界保持的尺度下考虑光滑试验函数 \(\varphi\) 的线性统计
\(
L_\kappa(\varphi)=\sum_{j=1}^{\kappa}\varphi(t_j)
\)。
则经中心化与方差归一化后 \(L_\kappa(\varphi)\) 收敛到标准正态分布，并且方差的极限由视界 Dirichlet--to--Neumann 生成元的半阶形式给出：
\[
\mathrm{Var}(L_\kappa(\varphi))
\ \to\
\frac{1}{2\pi}\int_{\RR}|\xi|\,|\widehat{\varphi_{\partial}}(\xi)|^2\,\dd\xi,
\]
其中 \(\varphi_{\partial}\) 为相应边界迹（或等价的调和投影）。
该猜想把涨落结构与命题 \ref{con:discussion-dtn-generator} 的 \(|D|\) 生成元在同一 \(H^{1/2}\) 半范数上严格对齐。
\end{conjecture}

\begin{corollary}[负谱体积的 Toeplitz 实现：\texorpdfstring{$\det\mathsf P$}{det P} 为有限证书负子空间正则化行列式的极限]\label{cor:xi-toeplitz-negative-regularized-determinant-limit}
\leanverified{paper\_xi\_toeplitz\_negative\_regularized\_determinant\_limit}
在稳定显影区间内，令 Toeplitz 证书 \(\mathbf{T}_N\) 具有 \(\kappa\) 个负特征值 \(\{\lambda_j^{-}(\mathbf{T}_N)\}_{j=1}^{\kappa}\)。
定义负子空间正则化行列式
\[
\mathrm{Det}_{-}(\mathbf{T}_N):=\prod_{j=1}^{\kappa}\bigl(-\lambda_j^{-}(\mathbf{T}_N)\bigr).
\]
则在半经典极限 \(N\to\infty\) 且 \(N(1-|a_j|)\to\infty\)（对全部缺陷点）下有
\[
\boxed{\;
\lim_{N\to\infty}\mathrm{Det}_{-}(\mathbf{T}_N)=\det\mathsf P(A).\;}
\]
\end{corollary}
\begin{proof}[证明要点]
由推论 \ref{cor:xi-toeplitz-negative-log-volume-limit} 有
\(\sum_{j=1}^{\kappa}\log(-\lambda_j^{-}(\mathbf{T}_N))\to \log\det\mathsf P\)；
两边取指数即得所示结论。证毕。
\end{proof}

\begin{corollary}[宏观分裂的 \texorpdfstring{$\kappa^2$}{k^2} 级行列式塌缩：交叉互能下界给出的指数压缩律（52）]\label{cor:xi-cross-energy-lower-bound-by-min-separation}
\leanverified{paper\_xi\_cross\_energy\_lower\_bound\_by\_min\_separation}
对并集 \(A\sqcup B\) 的交叉互能 \(\mathcal I(A;B)\)（命题 \ref{prop:xi-cross-energy-gaussian-mutual-information}），记
\(
\rho_{\min}(A,B):=\min_{a\in A,b\in B}\rho(a,b)
\)。
则有可审计下界
\[
\boxed{\;
\mathcal I(A;B)
=2\sum_{a\in A}\sum_{b\in B}\bigl(-\log\rho(a,b)\bigr)
\ \ge\
2\,|A|\,|B|\cdot\bigl(-\log\rho_{\min}(A,B)\bigr).\;}
\]
从而由 \(\mathcal I(A;B)\) 的定义得到严格压缩律
\[
\boxed{\;
\det\mathsf P(A\cup B)
=\det\mathsf P(A)\det\mathsf P(B)\,e^{-\mathcal I(A;B)}
\ \le\
\det\mathsf P(A)\det\mathsf P(B)\,\rho_{\min}(A,B)^{2|A||B|}\; }.
\]
特别地，若存在常数 \(\rho_0<1\) 使 \(\rho_{\min}(A,B)\le \rho_0\)，且 \(|A|\) 与 \(|B|\) 同阶增长，则相对塌缩比值满足
\(
\det\mathsf P(A\cup B)\big/(\det\mathsf P(A)\det\mathsf P(B))
\le \exp(-c\,|A||B|)
\),
其中 \(c=-2\log\rho_0>0\)，从而在宏观分裂 regime 中出现 \(\exp(-\Theta(\kappa^2))\) 级别的指数塌缩。
\end{corollary}
\begin{proof}[证明要点]
由命题 \ref{prop:xi-pick-poisson-union-submultiplicativity-cross-energy} 的恒等式
\(\mathcal I(A;B)=2\sum_{a\in A,b\in B}(-\log\rho(a,b))\)，并用
\(-\log\rho(a,b)\ge -\log\rho_{\min}(A,B)\) 即得第一条不等式。
将 \(\mathcal I(A;B)=-\log\det\mathsf P(A\cup B)+\log\det\mathsf P(A)+\log\det\mathsf P(B)\) 重写为
\(
\det\mathsf P(A\cup B)=\det\mathsf P(A)\det\mathsf P(B)e^{-\mathcal I(A;B)}
\)，
并代入 \(\mathcal I(A;B)\ge 2|A||B|(-\log\rho_{\min})\) 即得压缩律。证毕。
\end{proof}

\begin{theorem}[近极值锚定构型的宏观分裂阻断]\label{thm:xi-anchored-near-extremal-macroscopic-split-obstruction}
设有限锚定构型分解为不交并
\[
A=A^{(1)}\sqcup A^{(2)},
\qquad
n_i:=|A^{(i)}|,
\qquad
p_i:=\sum_{a\in A^{(i)}}P_a(-1)
\qquad(i=1,2).
\]
记
\[
\kappa:=n_1+n_2,
\qquad
p_{\Sigma}:=p_1+p_2,
\qquad
\varepsilon(A):=\kappa\log\!\Bigl(\frac{p_{\Sigma}}{\kappa}\Bigr)-\log\det\mathsf P(A).
\]
则有精确下界
\[
\boxed{\;
\varepsilon(A)\ge \mathcal I(A^{(1)};A^{(2)})
\ge
2n_1n_2\bigl(-\log\rho_{\min}(A^{(1)},A^{(2)})\bigr),\;}
\]
其中
\[
\rho_{\min}(A^{(1)},A^{(2)})
:=
\min_{a\in A^{(1)},\,b\in A^{(2)}}\rho(a,b).
\]
特别地，若存在常数 \(\eta>0\) 与 \(\rho_0<1\) 使
\[
n_1,n_2\ge \eta\kappa,
\qquad
\rho_{\min}(A^{(1)},A^{(2)})\le \rho_0,
\]
则
\[
\boxed{\;
\varepsilon(A)\ge 2\eta^2\kappa^2(-\log\rho_0).\;}
\]
因此任意满足 \(\varepsilon(A)=o(\kappa^2)\) 的近极值序列，都不能出现具有有界最小伪双曲分离度的宏观二分裂。
\end{theorem}
\begin{proof}
由交叉互能恒等式，
\[
\det\mathsf P(A)
=
\det\mathsf P(A^{(1)})\det\mathsf P(A^{(2)})\,e^{-\mathcal I(A^{(1)};A^{(2)})}.
\]
再由推论 \ref{cor:xi-anchored-capacity-extremal-fixed-flux}，对 \(i=1,2\) 有
\[
\det\mathsf P(A^{(i)})\le \Bigl(\frac{p_i}{n_i}\Bigr)^{n_i}.
\]
因此
\[
\log\det\mathsf P(A)
\le
n_1\log\!\Bigl(\frac{p_1}{n_1}\Bigr)
+
n_2\log\!\Bigl(\frac{p_2}{n_2}\Bigr)
-
\mathcal I(A^{(1)};A^{(2)}).
\]
对 \(n_1+n_2=\kappa\) 与 \(p_1+p_2=p_{\Sigma}\) 应用加权 AM--GM，得到
\[
\Bigl(\frac{p_1}{n_1}\Bigr)^{n_1}
\Bigl(\frac{p_2}{n_2}\Bigr)^{n_2}
\le
\Bigl(\frac{p_{\Sigma}}{\kappa}\Bigr)^\kappa,
\]
亦即
\[
n_1\log\!\Bigl(\frac{p_1}{n_1}\Bigr)
+
n_2\log\!\Bigl(\frac{p_2}{n_2}\Bigr)
\le
\kappa\log\!\Bigl(\frac{p_{\Sigma}}{\kappa}\Bigr).
\]
代回即得
\(
\varepsilon(A)\ge \mathcal I(A^{(1)};A^{(2)})
\)。
第二个下界正是推论 \ref{cor:xi-cross-energy-lower-bound-by-min-separation} 在 \((A^{(1)},A^{(2)})\) 上的应用。
若再有 \(n_i\ge \eta\kappa\) 且 \(\rho_{\min}\le \rho_0\)，则
\[
\varepsilon(A)
\ge
2n_1n_2(-\log\rho_0)
\ge
2\eta^2\kappa^2(-\log\rho_0).
\]
证毕。
\end{proof}

\begin{proposition}[并集合交叉互能的显式几何上界：以横向分离控制体积耦合]\label{prop:xi-pick-poisson-cross-energy-horizontal-separation-bound}
设缺陷多重集分解为不交并 \(A\bigsqcup B\subset\mathbb D\)。
沿用命题 \ref{prop:xi-pick-poisson-image-charge-green-energy} 的端点规一化上半平面坐标 \(t=\gamma+\mathrm{i}p\)（其中 \(p(t)=\Im t=P_a(-1)\)）。
若存在常数 \(L>0\) 使得对任意 \(a\in A\)、\(b\in B\) 都有
\[
|\Re t(a)-\Re t(b)|\ge L,
\]
则对任意交叉对 \((a,b)\in A\times B\) 成立
\[
1-\rho(a,b)^2
=\frac{4p(a)p(b)}{(\Re t(a)-\Re t(b))^2+(p(a)+p(b))^2}
\ \le\ \frac{4p(a)p(b)}{L^2}.
\]
若进一步满足
\(
\max_{a\in A,\ b\in B}\frac{4p(a)p(b)}{L^2}<1
\)，则交叉互能满足上界
\[
\mathcal I(A;B)
:=2\sum_{a\in A}\sum_{b\in B}\bigl(-\log\rho(a,b)\bigr)
=-\log\!\Bigl(\prod_{a\in A,\ b\in B}\rho(a,b)^2\Bigr)
\ \le\
\sum_{a\in A}\sum_{b\in B}
\frac{\tfrac{4p(a)p(b)}{L^2}}{1-\tfrac{4p(a)p(b)}{L^2}}.
\]
从而得到并集体积下界
\[
\det\mathsf P(A\cup B)
=\det\mathsf P(A)\det\mathsf P(B)\prod_{a\in A,\ b\in B}\rho(a,b)^2
\ \ge\
\det\mathsf P(A)\det\mathsf P(B)\prod_{a\in A,\ b\in B}\Bigl(1-\frac{4p(a)p(b)}{L^2}\Bigr).
\]
\end{proposition}
\begin{proof}
由命题 \ref{prop:xi-pick-poisson-image-charge-green-energy} 的半平面表示
\(
\rho(a,b)^2=\frac{|t(a)-t(b)|^2}{|t(a)-\overline{t(b)}|^2}
\)，
令 \(\Delta:=\Re t(a)-\Re t(b)\)，则
\[
|t(a)-t(b)|^2=\Delta^2+(p(a)-p(b))^2,\qquad
|t(a)-\overline{t(b)}|^2=\Delta^2+(p(a)+p(b))^2,
\]
从而
\[
1-\rho(a,b)^2
=\frac{|t(a)-\overline{t(b)}|^2-|t(a)-t(b)|^2}{|t(a)-\overline{t(b)}|^2}
=\frac{4p(a)p(b)}{\Delta^2+(p(a)+p(b))^2}.
\]
由 \(|\Delta|\ge L\) 得到第一条不等式。
再记 \(u_{ab}:=1-\rho(a,b)^2\in(0,1)\)。
则
\(
-\log\rho(a,b)^2=-\log(1-u_{ab})
\)。
当 \(u\in(0,1)\) 时有标准不等式 \(-\log(1-u)\le u/(1-u)\)，故
\[
\mathcal I(A;B)
=-\log\!\Bigl(\prod_{a,b}\rho(a,b)^2\Bigr)
=\sum_{a,b}-\log(1-u_{ab})
\le \sum_{a,b}\frac{u_{ab}}{1-u_{ab}}
\le \sum_{a,b}\frac{\tfrac{4p(a)p(b)}{L^2}}{1-\tfrac{4p(a)p(b)}{L^2}}.
\]
最后，由并集交叉互能恒等式（命题 \ref{prop:xi-pick-poisson-union-submultiplicativity-cross-energy}）
\(
\det\mathsf P(A\cup B)=\det\mathsf P(A)\det\mathsf P(B)\prod_{a,b}\rho(a,b)^2
\)，
并用 \(\rho(a,b)^2=1-u_{ab}\ge 1-\tfrac{4p(a)p(b)}{L^2}\) 得到体积下界。证毕。
\end{proof}

\begin{theorem}[远离碰撞时 Pick--Poisson 能量的双侧稳定性：以双曲距离给出可审计扰动界]\label{thm:xi-pick-poisson-energy-stability-away-from-collision}
设两组缺陷点配置 \(A=\{a_j\}_{j=1}^{\kappa}\)、\(A'=\{a'_j\}_{j=1}^{\kappa}\subset\mathbb D\)（按重数并带编号）。
令圆盘的双曲距离定义为
\[
d_{\mathbb D}(a,b):=2\,\operatorname{artanh}\rho(a,b)\qquad(a,b\in\mathbb D),
\]
并记
\[
\delta:=\max_{1\le j\le\kappa} d_{\mathbb D}(a_j,a'_j),
\qquad
d_{\min}:=\min_{1\le j<k\le\kappa} d_{\mathbb D}(a_j,a_k).
\]
若 \(d_{\min}>2\delta\)，则对 Pick--Poisson 能量
\(
S_{\mathrm{gen}}(A):=-\log\det\mathsf P(A)
\)
成立显式稳定性估计
\[
\boxed{\;
\bigl|S_{\mathrm{gen}}(A')-S_{\mathrm{gen}}(A)\bigr|
\ \le\
\kappa\,\delta
\;+\;
\frac{2\kappa(\kappa-1)\,\delta}{\sinh(d_{\min}-2\delta)}\;}
\]
从而得到体积扰动的双侧控制
\[
\exp\!\Bigl(-\kappa\delta-\frac{2\kappa(\kappa-1)\delta}{\sinh(d_{\min}-2\delta)}\Bigr)
\ \le\
\frac{\det\mathsf P(A')}{\det\mathsf P(A)}
\ \le\
\exp\!\Bigl(\kappa\delta+\frac{2\kappa(\kappa-1)\delta}{\sinh(d_{\min}-2\delta)}\Bigr).
\]
\end{theorem}
\begin{proof}
由行列式分解（命题 \ref{prop:xi-pick-poisson-det-factorization}；亦见定理 \ref{thm:xi-pick-poisson-vandermonde-resultant-identity}），有
\[
S_{\mathrm{gen}}(A)
=\sum_{j=1}^{\kappa}\bigl(-\log p(a_j)\bigr)
+2\sum_{1\le j<k\le\kappa}\bigl(-\log\rho(a_j,a_k)\bigr),
\qquad p(a):=P_a(-1).
\]
\par
\textbf{自能项。}
沿用命题 \ref{prop:xi-pick-poisson-image-charge-green-energy} 的端点规一化上半平面坐标 \(t(a)\in\mathbb H\)，则 \(p(a)=\Im t(a)\)，并且 \(d_{\mathbb D}(a,b)=d_{\mathbb H}(t(a),t(b))\)（M\"obius 等距）。
在上半平面模型中，函数 \(-\log\Im t\) 的双曲梯度范数恒为 \(1\)，从而
\[
\bigl|-\log p(a'_j)+\log p(a_j)\bigr|
=\bigl|-\log\Im t(a'_j)+\log\Im t(a_j)\bigr|
\le d_{\mathbb H}(t(a_j),t(a'_j))
=d_{\mathbb D}(a_j,a'_j)
\le \delta.
\]
对 \(j\) 求和得自能差的绝对值不超过 \(\kappa\delta\)。
\par
\textbf{互能项。}
对 \(j<k\) 记 \(d_{jk}:=d_{\mathbb D}(a_j,a_k)\)、\(d'_{jk}:=d_{\mathbb D}(a'_j,a'_k)\)。
由三角不等式，
\(
|d'_{jk}-d_{jk}|\le d_{\mathbb D}(a_j,a'_j)+d_{\mathbb D}(a_k,a'_k)\le 2\delta
\)，
并由假设 \(d_{\min}>2\delta\) 得 \(d'_{jk}\ge d_{\min}-2\delta>0\)。
令
\(
f(d):=-\log\rho=-\log\tanh(d/2)
\)，
则
\(
f'(d)=1/\sinh d
\)
且在 \((0,\infty)\) 上单调递减。
由均值定理，
\[
\bigl|f(d'_{jk})-f(d_{jk})\bigr|
\le \frac{|d'_{jk}-d_{jk}|}{\sinh(d_{\min}-2\delta)}
\le \frac{2\delta}{\sinh(d_{\min}-2\delta)}.
\]
对全部 \(\binom{\kappa}{2}\) 对求和并乘以能量分解中的系数 \(2\)，得到互能差的绝对值不超过
\(
\frac{2\kappa(\kappa-1)\delta}{\sinh(d_{\min}-2\delta)}
\)。
两部分相加即得盒中不等式。
最后对
\(
S_{\mathrm{gen}}(A')-S_{\mathrm{gen}}(A)=-\log\frac{\det\mathsf P(A')}{\det\mathsf P(A)}
\)
取指数即可得到行列式比值的双侧界。证毕。
\end{proof}

\endinput



\begin{corollary}[多块分裂的行列式指数塌缩：\texorpdfstring{$\kappa^2$}{k^2} 级压缩的全分区版本]\label{cor:xi-pick-poisson-det-multipartition-kappa2-collapse}
\leanverified{paper\_xi\_pick\_poisson\_det\_multipartition\_kappa2\_collapse}
设有限锚定集合 \(A\) 分解为互不交并集 \(A=\biguplus_{r=1}^{R}A_r\)，并记 \(\kappa_r:=|A_r|\)。
在命题 \ref{prop:xi-pick-poisson-det-factorization} 的行列式分解制度下，恒有精确分解
\[
\boxed{\;
\det\mathsf P(A)
=\Bigl(\prod_{r=1}^{R}\det\mathsf P(A_r)\Bigr)\,
\prod_{1\le r<s\le R}\ \prod_{a\in A_r,\ b\in A_s}\rho(a,b)^{2}\; }.
\]
从而得到统一上界：令
\(
\rho_{\min}(A_r,A_s):=\min_{a\in A_r,\ b\in A_s}\rho(a,b)
\)，则
\[
\boxed{\;
\det\mathsf P(A)\ \le\
\Bigl(\prod_{r=1}^{R}\det\mathsf P(A_r)\Bigr)\,
\prod_{1\le r<s\le R}\rho_{\min}(A_r,A_s)^{2\kappa_r\kappa_s}\; }.
\]
特别地，若存在常数 \(\rho_0<1\) 使对所有 \(r<s\) 有 \(\rho_{\min}(A_r,A_s)\le\rho_0\)，则
\[
\det\mathsf P(A)\ \le\ \Bigl(\prod_{r=1}^{R}\det\mathsf P(A_r)\Bigr)\,
\exp\!\Bigl(-2(-\log\rho_0)\sum_{1\le r<s\le R}\kappa_r\kappa_s\Bigr),
\]
呈现普适的 \(\kappa^2\)-级指数塌缩律。
\end{corollary}
\begin{proof}
将 \(\prod_{a<b}\rho(a,b)^2\) 按“块内对”与“跨块对”分解：块内对恰对应 \(\prod_r\det\mathsf P(A_r)/\prod_{a\in A}P_a(-1)\)，跨块对保持显式乘积，从而得到第一式。
第二式对跨块对逐项应用 \(\rho(a,b)\le \rho_{\min}(A_r,A_s)\) 立得；最后一式由取对数并整理得到。证毕。
\end{proof}

\begin{remark}[锚定容量的“加法--乘法”会计结构与两种复合方式]\label{rem:xi-anchored-capacity-accounting-structures}
记 \(\Phi(A):=-\log\det\mathsf P(A)=S_{\mathrm{gen}}(A)\)（定义 \ref{def:xi-generalized-defect-entropy-logdet}）。
则对并集 \(A\sqcup B\) 有严格会计恒等式
\[
\Phi(A\cup B)=\Phi(A)+\Phi(B)+\mathcal I(A;B),\qquad \mathcal I(A;B)\ge 0,
\]
其中 \(\mathcal I(A;B)\) 为交叉互能（命题 \ref{prop:xi-cross-energy-gaussian-mutual-information}）。
另一方面，在张量化（独立并行）复合下，指数型量 \(\det\mathsf P\) 乘法、而 \(\Phi\) 加法。
因此 \(\Phi\) 同时呈现两种结构：并集方向为严格超可加（含交互能项），并行方向为严格可加。
结合推论 \ref{cor:xi-anchored-capacity-stabilizer-covariance} 的稳定子群同调漂移，可把 \(\Phi\) 视为“视界数据结构”的规范记账势候选。
\end{remark}

\begin{conjecture}[连续缺陷的容量塌缩判别：连续毛发迫使离散锚定容量趋零]\label{conj:xi-continuous-hair-capacity-collapse}
令一般缺陷由 Herglotz/Clark 测度分解诱导，并设其包含非平凡奇异内因子（连续毛发）。
则对任意离散近似族 \(A_\kappa\subset\mathbb D\)（\(\kappa\to\infty\)）若其可见边界密度在局部意义下逼近该连续缺陷，则必有
\[
\mathrm{cap}_{-1}(A_\kappa)\ \longrightarrow\ 0.
\]
该猜想与推论 \ref{cor:xi-discrete-vs-continuous-hair-by-negative-inertia} 的“负惯性无界 \(\Leftrightarrow\) 连续毛发/无限点状自由度”相呼应：连续缺陷无法在固定锚定口径下以正容量方式离散化，其体积型可验证余量在极限上必被回收为零。
\end{conjecture}

\endinput


\endinput
